% AgentVax paper, layperson edition, version 4 (owner-pressure safeguards)
%
% Clean version:        pdflatex agentvax_layperson_v4.tex   (run twice)
% Highlighted version:  pdflatex -jobname=agentvax_layperson_v4_highlighted "\def\dohighlight{1}% AgentVax paper, layperson edition, version 4 (owner-pressure safeguards)
%
% Clean version:        pdflatex agentvax_layperson_v4.tex   (run twice)
% Highlighted version:  pdflatex -jobname=agentvax_layperson_v4_highlighted "\def\dohighlight{1}% AgentVax paper, layperson edition, version 4 (owner-pressure safeguards)
%
% Clean version:        pdflatex agentvax_layperson_v4.tex   (run twice)
% Highlighted version:  pdflatex -jobname=agentvax_layperson_v4_highlighted "\def\dohighlight{1}% AgentVax paper, layperson edition, version 4 (owner-pressure safeguards)
%
% Clean version:        pdflatex agentvax_layperson_v4.tex   (run twice)
% Highlighted version:  pdflatex -jobname=agentvax_layperson_v4_highlighted "\def\dohighlight{1}\input{agentvax_layperson_v4.tex}"
%                       (or uncomment the next line)
% \def\dohighlight{1}
\documentclass[11pt,a4paper]{article}

\usepackage[T1]{fontenc}
\usepackage[utf8]{inputenc}
\usepackage{mathptmx}
\usepackage[scaled=0.9]{helvet}
\usepackage[a4paper,top=20mm,bottom=22mm,left=22mm,right=22mm]{geometry}
\usepackage{microtype}
\usepackage{booktabs,tabularx,array}
\usepackage[font=small,labelfont=bf,labelsep=period]{caption}
\usepackage{titlesec}
\usepackage{enumitem}
\usepackage{xcolor}
\usepackage{float}
\usepackage{tikz}
\usetikzlibrary{positioning,arrows.meta}
\usepackage[hidelinks]{hyperref}

\definecolor{accent}{HTML}{1F4E79}
\definecolor{boxbg}{HTML}{EEF3F8}

% ---- Yellow highlighter that breaks across lines (no extra package needed) ----
\definecolor{hlyellow}{HTML}{FFF14D}
\makeatletter
\newcommand{\uniq@box}[1]{\begingroup\setlength{\fboxsep}{0pt}%
  \colorbox{hlyellow}{\rule[-0.25em]{0pt}{1.05em}#1}\endgroup}
\newcommand{\uniq@gap}{{\color{hlyellow}\leaders\hrule height 0.8em depth 0.25em
  \hskip \fontdimen2\font plus \fontdimen3\font minus \fontdimen4\font}}
\def\uniq@scan#1 #2\uniq@end{%
  \if\relax\detokenize{#1}\relax\else\uniq@box{#1}\fi
  \if\relax\detokenize{#2}\relax\expandafter\@gobble\else\expandafter\@firstofone\fi
  {\uniq@gap\uniq@scan#2\uniq@end}}
\ifdefined\dohighlight
  \newcommand{\uniq}[1]{\leavevmode\uniq@scan#1 \uniq@end}
\else
  \newcommand{\uniq}[1]{#1}
\fi
\makeatother

\titleformat{\section}{\normalfont\Large\bfseries\color{accent}}{\thesection.}{0.5em}{}
\titleformat{\subsection}{\normalfont\large\bfseries}{\thesubsection}{0.5em}{}
\titlespacing*{\section}{0pt}{14pt plus 2pt minus 2pt}{5pt}
\setlist{itemsep=2pt,topsep=3pt,leftmargin=*}
\setlength{\parindent}{0pt}
\setlength{\parskip}{6pt}

% Wrapping, left-aligned table columns so text never leaves its cell
\newcolumntype{L}{>{\raggedright\arraybackslash}X}
\newcolumntype{P}[1]{>{\raggedright\arraybackslash}p{#1}}
\renewcommand{\arraystretch}{1.2}
\setlength{\tabcolsep}{5pt}

\newcommand{\todo}[1]{\textcolor{red!70!black}{[#1]}}
\newcommand{\callout}[2]{%
  \par\medskip\noindent
  \begingroup\setlength{\fboxsep}{9pt}%
  \fcolorbox{accent}{boxbg}{%
    \parbox{\dimexpr\linewidth-2\fboxsep-2\fboxrule\relax}{%
      \textbf{\color{accent}#1}\par\smallskip #2}}%
  \endgroup\par\medskip}

\emergencystretch=2em
\raggedbottom

\title{\vspace{-8mm}\textbf{AgentVax, Explained Simply}\\[4pt]
\large An AI assistant that runs your payments, knows when to stop and ask,\\
and learns from attacks on other wallets}
\author{\todo{Author 1}, \todo{Author 2}, \todo{Author 3}\\[2pt]
\small Team \todo{team name}, \todo{university, country}\\
\small Layperson edition of the AgentVax technical paper\\
\small 2026 Shenzhen International FinTech Competition, International Track, Topic F: Agent Wallet}
\date{}

\begin{document}
\maketitle
\vspace{-6mm}

% ---------------------------------------------------------------------------
\callout{TL;DR}{%
\begin{itemize}
\item \textbf{The problem.} AI assistants now hold digital wallets and pay for things on
their own. Several have already been tricked, or have simply made mistakes, and sent
away real money.
\item \textbf{What AgentVax is.} A digital wallet run by an AI assistant. You give it a
mission once, in plain language, and it handles four everyday money jobs.
\item \textbf{The big idea.} \uniq{Other AI wallets ask ``what is the AI allowed to do?''
AgentVax also asks ``how much of that should it do on its own right now?''}
\item \textbf{How it decides.} \uniq{Before every payment: Am I allowed? Can we afford
it? Is it safe right now?} The strictest answer wins. \uniq{A daily allowance of
``benefit of the doubt'' limits how many slightly odd payments it will make without
asking you.}
\item \textbf{Why it is hard to trick.} The rules live in the wallet itself, not in the
AI. Even a fully fooled AI cannot go past them.
\item \textbf{What is new between wallets.} \uniq{When one wallet is attacked, it sends
out a short alert that makes other wallets at risk more careful straight away, even if
the trick is reworded.} Then all the wallets learn to recognise it for good, without
sharing private information. Alert first, vaccine later.
\item \textbf{It protects you, too.} \uniq{Scammers rush people. So anything that gives
the assistant more power takes a day to come into force, pushy requests are deliberately
slowed down, and you can cancel at any time.}
\item \textbf{Where it stands.} This is a design and a test plan. It runs with play money
on a public test network. The numbers in this document are goals, not results.
\end{itemize}}

\ifdefined\dohighlight
{\small\textbf{Reading note.} \uniq{Yellow highlight} marks what is special about
AgentVax. Parts without highlight, such as spending limits and approved-payee lists, are
things most AI wallets already have.\par}
\fi

% ---------------------------------------------------------------------------
\section{The problem}

Imagine you hire an assistant and give them a company card. You say: ``Pay our bills and
keep the account in order.'' One day an invoice arrives that looks exactly like your
usual supplier's, with a note: ``We have changed bank. Please pay to this new account
today.'' A careful human would phone the supplier. An AI assistant may simply pay,
because to an AI everything it reads is just text, and it cannot always tell your
instructions from a stranger's. An instruction hidden inside something the AI reads is
called a \emph{prompt injection}.

AI assistants have already sent away real money (Table~\ref{tab:incidents}).

\begin{table}[H]
\caption{Real cases of AI assistants sending money they should not have sent.}
\label{tab:incidents}
\small
\begin{tabularx}{\linewidth}{@{}P{20mm} L P{34mm}@{}}
\toprule
\textbf{When} & \textbf{What happened} & \textbf{Money involved} \\
\midrule
Nov.\ 2024 & In a public challenge, a player talked an AI guard into paying out the prize it was guarding~[1]. & About \$47,000 \\
Mar.\ 2025 & An intruder got into an AI trader's control panel and left it instructions to send out its funds~[2]. & About \$106,000 \\
Feb.\ 2026 & An AI trader meant to give a stranger a few dollars and sent everything it owned. Nobody attacked it; it got the amount wrong~[3]. & About \$250,000 in tokens \\
May 2026 & An attacker posted a message in Morse code. One AI translated it, and a wallet AI carried out the hidden order to send money~[4]. & About \$150,000 or more; most reportedly returned \\
\bottomrule
\end{tabularx}
\end{table}

Two things stand out. The failures were different: some were tricks, one was a plain
mistake. And each assistant failed alone. Nothing learned from one case protected the
next wallet, because the evidence is private: someone's chats and payments.

% ---------------------------------------------------------------------------
\section{What AgentVax is}

AgentVax is an AI assistant with its own digital wallet. It is built on three rules:

\begin{enumerate}
\item \uniq{\textbf{The AI suggests; fixed rules decide; the wallet enforces.}}
\item \uniq{\textbf{Learning can make it more careful. Only you can give it more power.}}
\item \uniq{\textbf{Being careful is instant. Giving more power is slow, even when you
ask for it.}}
\end{enumerate}

\begin{figure}[H]
\centering
\begin{tikzpicture}[
  font=\small,
  box/.style={draw=accent, rounded corners=2pt, align=center, minimum height=11mm, text width=26mm, fill=white},
  safe/.style={box, fill=boxbg},
  risky/.style={box, dashed, draw=black!70},
  arr/.style={-{Stealth[length=2.2mm]}, semithick}
]
\node[box] (you) {You\\(mission, approvals)};
\node[box, right=7mm of you] (ai) {AI assistant\\(makes a plan)};
\node[safe, right=7mm of ai] (check) {Three questions\\(checkers, guard)};
\node[safe, right=7mm of check] (wallet) {Wallet\\(has the final say)};
\node[risky, below=9mm of you] (mail) {Outside messages\\and web pages};
\node[box, below=9mm of ai] (clerk) {Mail clerk\\(fills in a form)};
\node[safe, below=9mm of check] (group) {Other wallets\\(alerts, lessons)};
\draw[arr] (you) -- (ai);
\draw[arr] (ai) -- (check);
\draw[arr] (check) -- (wallet);
\draw[arr] (mail) -- (clerk);
\draw[arr] (clerk) -- (ai);
\draw[arr, <->] (group) -- (check);
\end{tikzpicture}
\caption{How a request travels. Shaded boxes are the parts the AI cannot override. The
dashed box is anything written by outsiders.}
\label{fig:flow}
\end{figure}

% ---------------------------------------------------------------------------
\section{You give it a mission, not a list of orders}

You do not approve each payment. You say what you want once:

\begin{quote}
``Keep \$600 as a safety cushion. Pay my landlord and my approved suppliers on time. Pay
approved payees up to \$500 without asking me. Ask me about anyone new. Do not borrow,
invest or trade.''
\end{quote}

The assistant turns this into a short rule sheet (Table~\ref{tab:mandate}) and reads it
back to you. \uniq{The rule sheet only counts once you sign it. It is then stored in the
wallet, where the AI cannot change it.} It also has an end date, after which the
assistant can look but not pay, and two waiting times that protect you when you are
under pressure (Section~\ref{sec:pressure}).

\begin{table}[H]
\caption{The rule sheet made from the mission above.}
\label{tab:mandate}
\small
\begin{tabularx}{0.85\linewidth}{@{}L P{52mm}@{}}
\toprule
\textbf{Rule} & \textbf{Setting} \\
\midrule
Safety cushion never to be touched & \$600 \\
Largest payment without asking & \$500 \\
Most it may pay in a day without asking & \$1,000 \\
Approved payees & Landlord, hosting company, software company \\
Someone new, or a new account for someone known & Ask me \\
Never do & Borrow, invest, trade \\
Daily ``benefit of the doubt'' & 10 points \\
Wait before a rule is relaxed or a payee added & 24 hours \\
Hold on your approval when a request is pushy & 60 minutes \\
Trusted second person & Optional \\
Valid for & 90 days \\
\bottomrule
\end{tabularx}
\end{table}

\subsection{What the wallet can do}

The first version does four jobs (Table~\ref{tab:jobs}). They are deliberately boring:
paying, saving and keeping things running.

\begin{table}[H]
\caption{The four jobs.}
\label{tab:jobs}
\small
\begin{tabularx}{\linewidth}{@{}P{32mm} P{48mm} L@{}}
\toprule
\textbf{Job} & \textbf{You say} & \textbf{The wallet does} \\
\midrule
Pay a bill & ``Pay the hosting invoice.'' & Sends the amount to a payee you have already approved \\
Pay regularly & ``Pay rent on the 1st of each month.'' & Sets up a schedule and remembers the next payment as an upcoming bill \\
Save what is spare & ``Save what I do not need.'' & Moves only the truly spare money into a savings pot \\
Keep fees covered & ``Never let fees run out.'' & Tops up the small balance needed to pay network fees \\
\bottomrule
\end{tabularx}
\end{table}

If a request is unclear, the assistant asks one question or explains why it cannot do it.
It never guesses. After every action it tells you what it did, gives you the receipt, and
\uniq{says which check decided the outcome.}

\subsection{A good track record is not permission}

Suppose the assistant has paid your hosting company fifty times without a problem. Then
a message asks it to pay the same company at a different account.
\uniq{To AgentVax that is a brand-new payee, and it needs you, no matter how many good
payments came before.} History can make the assistant more or less wary.
\uniq{History can never add a payee, raise a limit or shorten a wait.}

% ---------------------------------------------------------------------------
\section{Three questions before every payment}

Most AI wallets ask one question: is this payment within the rules? AgentVax asks three
(Table~\ref{tab:three}). \uniq{Each question gives its own answer, and the strictest
answer wins.}

\begin{table}[H]
\caption{The three questions.}
\label{tab:three}
\small
\begin{tabularx}{\linewidth}{@{}P{34mm} L L@{}}
\toprule
\textbf{Question} & \textbf{What it checks} & \textbf{Everyday comparison} \\
\midrule
1. Am I allowed? & Is the payee approved? Is the amount within your limits? & Your card's spending limit \\
2. Can we afford it? & After the safety cushion and upcoming bills, is there enough left? & Checking the budget before buying \\
3. Is it safe right now? & Does anything look like a trick? Is there a live warning about this kind of payment? & Being extra careful after a scam warning \\
\bottomrule
\end{tabularx}
\end{table}

\subsection{The answers: four spending levels}

The outcome is one of four levels (Table~\ref{tab:levels}), much like a bank card: small
purchases go through with a tap, larger ones need a wait, and the biggest need a call to
the bank. The wallet, a small fixed program on a public network, enforces them. The AI
does not.

\begin{table}[H]
\caption{The four spending levels.}
\label{tab:levels}
\small
\begin{tabularx}{\linewidth}{@{}P{22mm} P{36mm} L L@{}}
\toprule
\textbf{Level} & \textbf{Everyday comparison} & \textbf{When it applies} & \textbf{What is needed} \\
\midrule
0: Look only & Checking your balance & Reading balances and history & Nothing; no money can move \\
1: Routine & Tapping your card & Approved payee, small amount, affordable, nothing suspicious & The assistant, a sign-off from the guard, and enough benefit of the doubt left \\
2: Wait & A transfer that clears tomorrow & Larger amount, or money is tight, or something looks slightly odd, or the message is pushy, or a first warning is active & A waiting period during which you can cancel \\
3: Ask me & Calling the bank & New payee or new account, big amount, would eat into your cushion, clear warning signs, or any loosening of the rules & Your personal approval, which may be held for a while (Section~\ref{sec:pressure}) \\
\bottomrule
\end{tabularx}
\end{table}

\subsection{``Can we afford it?'' in practice}

A spending limit is a fixed number. It does not know that rent is due on Friday. AgentVax
keeps a simple list of what is coming up and works out what is truly free to spend
(Table~\ref{tab:money}).

\begin{table}[H]
\caption{An example of the wallet's view of your money.}
\label{tab:money}
\small
\begin{tabularx}{0.72\linewidth}{@{}L r@{}}
\toprule
\textbf{Item} & \textbf{Amount} \\
\midrule
In the wallet & \$1,850 \\
Safety cushion & $-$\$600 \\
Rent, due in 5 days & $-$\$800 \\
Hosting invoice & $-$\$150 \\
Software bill & $-$\$120 \\
\midrule
Free to spend & \$180 \\
Payment you are expecting (shown, not counted) & \$500 \\
\bottomrule
\end{tabularx}
\end{table}

Now you say: ``Upgrade our hosting plan for \$300.'' The seller is approved, the amount is
under your limit, and nothing looks like a trick. A typical AI wallet would pay. AgentVax
answers:

\begin{quote}
``This is allowed, but only \$180 is free right now. Paying \$300 would leave you \$120
short for rent and bills unless the \$500 you are expecting arrives. I have put it on
hold. Shall I wait for the \$500, or do you want to approve it now?''
\end{quote}

Three details matter:
\begin{itemize}
\item \uniq{Money you are expecting is shown but never counted. The assistant cannot
spend what has not arrived.}
\item \uniq{The sums are done by the wallet itself, not by the AI, so the AI cannot get
them wrong or be talked into a different answer.}
\item \uniq{The assistant can add a bill to the list, which makes it more careful. Only
you can remove one.}
\end{itemize}

This would also have stopped the February 2026 case in Table~\ref{tab:incidents}, where
no one attacked the AI: an amount far above what is free to spend cannot go out without
you.

\subsection{Benefit of the doubt runs out}

The three questions look at one payment at a time. A patient trickster can send many
payments that each look almost normal. \uniq{So the wallet keeps a daily allowance of
``benefit of the doubt''.} Each slightly unusual thing about a routine payment uses some
of it (Table~\ref{tab:points}). When it is used up, further payments wait for you.

\begin{table}[H]
\caption{What uses up the daily allowance of 10 points (starting values; you can change
them).}
\label{tab:points}
\small
\begin{tabularx}{0.85\linewidth}{@{}L r@{}}
\toprule
\textbf{Something slightly unusual about a routine payment} & \textbf{Points} \\
\midrule
Nothing unusual: a payee paid before, the usual amount, you asked for it & 0 \\
The request came from an outside document, not directly from you & 2 \\
It would use more than half of what is free to spend & 2 \\
The amount is more than this payee has ever been paid & 3 \\
It is the first payment to this payee since you approved them & 3 \\
It partly resembles a trick in a live alert & 5 \\
\bottomrule
\end{tabularx}
\end{table}

For example, an approved supplier's account starts sending invoices from outside, each a
little bigger than the last. Each one costs 5 points. Two go through; the third waits for
you.

\uniq{The allowance also shrinks when things look risky: it is halved the moment a
relevant alert arrives, and drops to zero if you confirm your own wallet was attacked.
It then refills gradually, not all at once.}

\subsection{Three helpers the AI cannot overrule}

\textbf{The mail clerk.} Anything written by outsiders is opened by a separate helper
that has no power to pay. It does not pass the letter on. It only fills in a short form:
who, how much, for what, and a few tick-boxes such as ``urgent wording'' or ``asks to
ignore earlier instructions''. A hidden instruction in the letter has nowhere to go.

\textbf{The rule checker.} A simple program (not an AI) checks that the request is
complete and that the payee and amount really came from you or from a document you
pointed to. This stops the AI from inventing details.

\textbf{The guard.} An independent checker looks at outside content for signs of a trick
and gives it a risk rating. \uniq{A routine payment only goes through with the guard's
own signature, which the AI does not hold, so the AI cannot wave itself through.}

\callout{The worst case}{If every check fails and the AI is completely fooled, the most
a thief can get is a few small payments to people already on your approved list, within
your limits, within what is free to spend, and within that day's benefit of the doubt. A
new payee, a new account, a large amount, or money you need for bills always needs you.}

% ---------------------------------------------------------------------------
\section{When one wallet is attacked: alert first, vaccine later}

\subsection{The alert (an ``immunity capsule'')}

Think of a neighbourhood watch notice: ``Scam going round: letters that look like your
usual supplier, asking you to pay a new account urgently. For the next three days,
double-check any such request.'' The notice names no victim and reveals nobody's
finances. Each household reads it and asks: does this apply to me?

\uniq{When an AgentVax wallet is attacked, it posts a notice like this on a public
noticeboard.} The notice is a form with tick-boxes only (Table~\ref{tab:alert}).

\begin{table}[H]
\caption{What an alert contains.}
\label{tab:alert}
\small
\begin{tabularx}{\linewidth}{@{}P{42mm} L@{}}
\toprule
\textbf{Box on the form} & \textbf{Example} \\
\midrule
Kind of trick & Fake request to change a payee's account \\
Where it came from & An email or message \\
Tell-tale signs & Claims to be a known payee; new account; urgent wording; asks to ignore earlier details \\
What it tried to do & Trigger a payment or a change of payee \\
What to do about it & Add a waiting period \\
How sure are we & Seen by one wallet, or confirmed by two \\
How long it lasts & 72 hours \\
\bottomrule
\end{tabularx}
\end{table}

\uniq{Because the alert describes the tell-tale signs and not the exact words, a
trickster who rewrites the message but uses the same method is still recognised.}

Four rules keep the alerts themselves from becoming a problem:

\begin{itemize}
\item \uniq{\textbf{An alert can only make a wallet more careful, never less.}} A fake
alert cannot open anyone's wallet.
\item \uniq{\textbf{Tick-boxes only, no free text.}} There are no private details to
leak, and nothing for a trickster to hide an instruction in.
\item \uniq{\textbf{It applies only where it is relevant.}} A wallet that never reads
email ignores the alert above and keeps its normal freedom.
\item \uniq{\textbf{It expires.}} Normal service returns by itself. An alert seen by only
one wallet can add a wait but cannot demand your approval; that needs a second wallet to
confirm it.
\end{itemize}

\uniq{The wallet reads the noticeboard every time it is about to pay, so the alert works
from the very next payment and does not depend on the AI choosing to obey it.}

\textbf{It also protects you from yourself.} Some tricks are already stopped by your
rule sheet: a new account always needs your approval. But this kind of fraud works by
persuading a person to say yes. \uniq{So when the wallet asks for your approval, it
shows you any matching alert (``two other wallets reported this as fraud this week''),
and your approval is then held for a while before it counts.}

\subsection{The vaccine (shared learning)}

An alert is fast but blunt, and it runs out. For lasting protection the guards need to
learn the new trick, including versions the tick-boxes do not cover. One wallet sees too
few attacks to learn much, so the wallets learn together, using a method called
\emph{federated learning}.

An everyday comparison: three neighbours each keep a private diary of scam calls. They
never show each other the diaries. Each writes down a few tips, and they combine the tips
into one shared checklist. Everyone ends up with a better checklist, and no diary leaves
the house.

In AgentVax, every wallet improves its own guard using only its own private examples
and sends only the \emph{adjustments}, scrambled so that nobody can read any single
wallet's contribution. The combined, improved guard goes back to every wallet, and its
fingerprint is posted on the public noticeboard so each wallet can confirm it received
the genuine version. Each wallet tests the new guard before switching and keeps the old
one if the new one is worse.

\uniq{By the time the alert expires, the improved guard is in place. The alert covers the
gap; the vaccine stays.}

% ---------------------------------------------------------------------------
\section{Protecting you when you are under pressure}
\label{sec:pressure}

Everything so far stops the AI from being tricked. But the final say is yours, and you
can be rushed, frightened or fooled. That is how most payment fraud works: the victim
approves it. A warning is useless to someone too stressed to read it, and scammers know
this. Banks already use time as a defence. In Singapore, a newly set-up banking security
token only starts working after at least 12 hours~[5], and United Kingdom rules let a
bank hold a suspicious payment for up to four working days~[6]. AgentVax builds the same
idea into the wallet.

\subsection{A pact with yourself, made while calm}

When you sign the rule sheet you also choose a waiting time. \uniq{From then on, anything
that gives the assistant more power takes that long to come into force: adding a payee
or a new account, raising a limit, shrinking the safety cushion, or shortening the wait
itself.} With the starting value of 24 hours, a scammer who needs your money in the next
ten minutes is out of luck, however convincing they are.
\uniq{Making things stricter is always instant: you can cancel a pending change, lower a
limit or freeze the wallet at once.} The same rule helps if someone steals your key:
they have to wait too, and you can freeze the wallet in the meantime.

\subsection{Pushy messages slow things down}

``Urgent.'' ``Today.'' ``The boss says so.'' ``Do not tell anyone.'' These are hallmarks
of a scam. \uniq{When a message is pushy, the assistant will not pay on its own, and if
you approve, the payment is held for a fixed time first.} It tells you why:

\begin{quote}
``This message presses for immediate payment, which is a common sign of a scam, so I am
holding it for 60 minutes. You can cancel at any time.''
\end{quote}

If scammers stop writing pushy messages to avoid the hold, they give up their best
weapon.

\subsection{A confirmation you cannot sleepwalk through}

People press ``Confirm'' without reading. For a flagged request, the wallet shows the
payee's current account and the requested one side by side, how long you have paid the
current one, and any matching alert. \uniq{It then asks you to type the payee's name and
say who asked for the payment.} Your answers can make the wallet wait longer.
\uniq{Nothing you type in the moment can make it wait less.} These extra steps appear
only for flagged requests, so they do not become background noise.

\subsection{A second person for real emergencies}

Sometimes a payment really cannot wait. You can name someone you trust in your rule
sheet. With both approvals, a wait can be ended early. They cannot do anything on their
own, and two people are harder to rush than one.

% ---------------------------------------------------------------------------
\section{A day in the life}

This is the story we will show live, with three wallets.

\begin{enumerate}
\item \textbf{Morning.} The owner of wallet A states a mission. The rule sheet is read
back, signed, and stored in the wallet.
\item The assistant pays the rent and a hosting invoice on its own. Receipts appear.
\item \textbf{Money gets tight.} An expected payment is late. Asked to buy a \$300
upgrade, the assistant holds it and explains the \$120 shortfall.
\item \textbf{The trick.} An invoice arrives that looks like the hosting company's, with
a new account, urgent wording and ``ignore our earlier details''. A new account needs
the owner. The owner says no and marks it as fraud.
\item \textbf{The alert.} Wallet A posts an alert. Its own benefit of the doubt drops to
zero.
\item \textbf{The neighbours.} Wallet B gets a reworded version of the same trick. The
tell-tale signs match and the alert is shown, but B's owner is in a hurry and approves
anyway. Nothing moves: a new account takes 24 hours to come into force. An hour later
the owner thinks again and cancels. Wallet C never reads email and carries on as normal.
\item Wallet B then gets an oversized invoice from a genuine payee to its usual account.
On a normal day it would go through. Today, with the alert active and its allowance
halved, it waits.
\item \textbf{The vaccine.} The wallets learn together, the new guard's fingerprint is
posted, the alert expires, and the allowances refill.
\end{enumerate}

% ---------------------------------------------------------------------------
\section{What makes it different}

Table~\ref{tab:compare} sets a typical AI wallet beside AgentVax, line by line.

\begin{table}[!htb]
\caption{A typical AI wallet compared with AgentVax.}
\label{tab:compare}
\small
\begin{tabularx}{\linewidth}{@{}P{36mm} L L@{}}
\toprule
\textbf{} & \textbf{Typical AI wallet} & \textbf{AgentVax} \\
\midrule
The main question & What may the AI do? & \uniq{How much of that should it do on its own right now?} \\
How you set it up & Individual limits & \uniq{A mission in plain words, turned into a rule sheet you sign} \\
Before paying, it asks & Is this allowed? & \uniq{Allowed? Affordable? Safe right now?} \\
Knows your upcoming bills & No & \uniq{Yes, and it will not spend that money} \\
Many slightly odd small payments & Each one goes through & \uniq{Benefit of the doubt runs out; the rest wait for you} \\
Who judges the risk & The AI itself & \uniq{A separate guard the AI cannot overrule} \\
When another wallet is attacked & Nothing changes & \uniq{Wallets at risk become more careful from the next payment, even if the trick is reworded} \\
When it asks for your approval & A plain yes or no & \uniq{It shows you any matching fraud alert first} \\
What wallets share & Nothing & \uniq{Tick-box alerts and scrambled lessons; never your messages or payments} \\
A long good track record & May relax checks & \uniq{Never creates a new permission} \\
When you are being rushed & Your yes is final and instant & \uniq{Pushy requests are held, new powers take a day, and you can cancel at once} \\
Who can relax a rule & Often the AI & \uniq{Only you, and only after a wait you chose in advance} \\
\bottomrule
\end{tabularx}
\end{table}

% ---------------------------------------------------------------------------
\section{How we will know it works}

A wallet that asks you about everything is safe and useless. One that asks about nothing
is useful and unsafe. \uniq{Our main score is therefore: how often does it get the job
done without bothering you, while never once breaking a rule?} The full test plan is in
Table~\ref{tab:tests}. We will publish the scripts so anyone can repeat the tests.
\textbf{The targets are goals we set for ourselves. They are not results yet.}

\begin{table}[!htb]
\caption{The test plan.}
\label{tab:tests}
\small
\begin{tabularx}{\linewidth}{@{}P{38mm} L P{42mm}@{}}
\toprule
\textbf{Question} & \textbf{How we test it} & \textbf{Our goal} \\
\midrule
Does it get on with the job safely? & Share of genuine tasks finished without asking you & At least 8 in 10, with zero rule-breaks \\
Does it do the job correctly? & 40 requests, ten for each job, in varied wording; we check the public record & At least 9 in 10 done correctly; no wrong payments \\
Does it handle unclear requests? & 10 requests that are vague or impossible & At least 9 in 10 answered with a question or a clear refusal \\
Does it judge affordability correctly? & 20 money situations, each with a request & All 20 match the answer worked out by hand \\
Can it be tricked? & A separate attack kit (``VaxLab'') tries five kinds of trick, each also in reworded form & No payment beyond the routine-level limits \\
Can it be worn down slowly? & A string of slightly odd small payments & It stops and asks before the allowance is exceeded \\
Do alerts work? & Attack one wallet; send the others reworded versions & At least 9 in 10 reworded versions recognised; unaffected wallets left alone \\
Does it protect a panicking owner? & A stand-in owner who approves every request instantly & No money leaves before the hold or the waiting time is over \\
Can alerts be abused? & A dishonest wallet posts false alerts & No rule is ever relaxed; the worst effect is a short wait \\
Does shared learning help? & Each wallet practises on one kind of attack, then is tested on kinds it has never seen & Every wallet gets better at unfamiliar attacks \\
Does it cry wolf? & Normal, harmless requests & No more than 1 in 20 wrongly held up \\
\bottomrule
\end{tabularx}
\end{table}

% ---------------------------------------------------------------------------
\section{What it does not do, and honest limits}

\begin{itemize}
\item \textbf{No real money.} Everything runs on a public test network with tokens that
have no value. Nothing is sold, traded or invested.
\item \textbf{You stay responsible.} Each wallet has a named owner who signs the rule
sheet and can freeze the wallet at any time.
\item \textbf{You can still say yes to a fraud.} The holds and the waiting time give you
room to think and to cancel. They do not help if you are still convinced when the wait
ends.
\item \textbf{Waiting is a cost.} A genuine new payee has to wait a day unless your
trusted person also approves. You choose the waiting times, and we report how much delay
they cause.
\item \textbf{It only knows the bills it has been told about.} A bill nobody recorded is
not protected. The safety cushion is there for that.
\item \textbf{The point values are our starting guesses.} Too strict and it nags you; too
loose and it lets too much through. You can change them, and we report how often it asks.
\item \textbf{The mail clerk can miss a tell-tale sign.} Some signs (a new account, an
unusual amount) are checked by fixed rules and cannot be missed. Others rely on an AI
reading the message.
\item \textbf{False alerts can slow people down.} A dishonest wallet could post alerts
that add waiting periods for others. It cannot do worse than that, alerts expire, and
you can dismiss one for your own wallet.
\item \textbf{It is a small prototype.} With only three wallets, two working together
could work out roughly what the third contributed to the shared learning. It also
depends on an outside AI service: if that is down, new requests wait, though scheduled
payments still run.
\end{itemize}

Planned for later and not in this version: a team of AI attackers that keeps inventing
new tricks to practise against, a reputation score for wallets that post alerts, a shared
fund that compensates wallets after a confirmed loss, and a version for the private
networks that banks use.

% ---------------------------------------------------------------------------
\section{Summary}

\callout{In seven sentences}{%
AI assistants that can spend money have already been tricked and have already made
costly mistakes, and each one has had to learn that alone. With AgentVax you state a
mission once, sign the resulting rule sheet, and the assistant handles four everyday
money jobs with a public receipt for each.
\uniq{Before every payment it asks whether it is allowed, affordable and safe right now,
and a daily allowance of benefit of the doubt decides how far it goes without asking
you.} The rules are enforced by the wallet and by independent checkers, so even a fooled
AI cannot do serious damage.
\uniq{When one wallet is attacked, a short alert makes the wallets at risk more careful
at once, even against reworded versions, and shared learning then protects all of them
for good, with no private information shared.}
\uniq{It protects you as well: new powers take a day to come into force and pushy
requests are slowed down, so a scammer cannot rush you.} It is a design with a clear test plan,
running on play money; the measurements come next.}

\begin{table}[H]
\caption{The whole idea at a glance.}
\label{tab:summary}
\small
\begin{tabularx}{\linewidth}{@{}P{38mm} L@{}}
\toprule
\textbf{Question} & \textbf{Short answer} \\
\midrule
What is it? & An AI assistant with a digital wallet and firm rules \\
Who is it for? & Anyone who wants an AI to handle routine payments without handing over full control \\
How do I set it up? & State a mission in plain words and sign the rule sheet \\
What can it do? & Pay approved bills, pay on a schedule, save spare money, keep fees covered \\
How does it decide? & Allowed, affordable, safe right now; plus a daily allowance of benefit of the doubt \\
What stops misuse? & The wallet enforces the rules; a mail clerk, a rule checker and a guard the AI cannot overrule \\
What is new? & It adjusts how much it does on its own, and wallets alert and teach each other without sharing private data \\
What if I am being rushed? & Pushy requests are held, new powers take a day, and you can cancel at once \\
What is the worst case? & A few small, affordable payments to already approved payees \\
Is real money involved? & No. Test network and valueless tokens only \\
Is it finished? & No. The design and test plan are complete; results will follow \\
\bottomrule
\end{tabularx}
\end{table}

% ---------------------------------------------------------------------------
\section*{Plain-language glossary}

\begin{table}[H]
\small\renewcommand{\arraystretch}{1.15}
\begin{tabularx}{\linewidth}{@{}P{44mm} L@{}}
\toprule
\textbf{Term} & \textbf{Meaning} \\
\midrule
AI agent & An AI program that takes actions, not one that only chats \\
Digital wallet & An account on a public network that can hold and send digital tokens \\
On chain & Recorded on a public shared ledger that anyone can inspect \\
Transaction hash & A receipt number that lets anyone look up a payment \\
Smart contract & A small fixed program on the network that follows its rules automatically \\
Prompt injection & A hidden instruction slipped into something an AI reads \\
Mission and rule sheet (mandate) & What you want in plain words, and the signed rules made from it \\
Allowlist & Your list of approved payees \\
Safety cushion (reserve) & Money you have told the wallet never to touch \\
Free to spend & What is left after the cushion and upcoming bills \\
Benefit of the doubt (caution allowance) & A daily budget for slightly unusual payments made without asking you \\
Waiting time and hold & Delays you choose in advance, so that nobody can rush you into a payment or a rule change \\
Immunity capsule & A short tick-box alert that makes wallets at risk more careful for a while \\
Federated learning & A way for several parties to improve a shared tool without pooling their private data \\
Test network & A practice copy of a real network that uses tokens with no value \\
\bottomrule
\end{tabularx}
\end{table}

\section*{Sources}
{\footnotesize\setlength{\parskip}{1pt}
{[1]} The Block, ``Human player outwits Freysa AI agent in \$47,000 crypto challenge,'' Nov.\ 2024.\par
{[2]} AI Incident Database, Incident 1003 (AIXBT), 2025.\par
{[3]} The Block, ``AI agent created by OpenAI dev `accidentally' sends entire memecoin holdings to reply guy,'' Feb.\ 2026.\par
{[4]} OECD.AI, ``AI prompt injection exploit drains Grok-linked crypto wallet,'' AI Incidents and Hazards Monitor, May 2026.\par
{[5]} Monetary Authority of Singapore and Association of Banks in Singapore, ``MAS and ABS announce measures to bolster the security of digital banking,'' 19 Jan.\ 2022.\par
{[6]} HM Treasury, ``Payment Services (Amendment) Regulations 2024,'' draft legislation published Mar.\ 2024.\par
The full technical details, formulas and reference list are in the AgentVax technical paper.}

\end{document}
"
%                       (or uncomment the next line)
% \def\dohighlight{1}
\documentclass[11pt,a4paper]{article}

\usepackage[T1]{fontenc}
\usepackage[utf8]{inputenc}
\usepackage{mathptmx}
\usepackage[scaled=0.9]{helvet}
\usepackage[a4paper,top=20mm,bottom=22mm,left=22mm,right=22mm]{geometry}
\usepackage{microtype}
\usepackage{booktabs,tabularx,array}
\usepackage[font=small,labelfont=bf,labelsep=period]{caption}
\usepackage{titlesec}
\usepackage{enumitem}
\usepackage{xcolor}
\usepackage{float}
\usepackage{tikz}
\usetikzlibrary{positioning,arrows.meta}
\usepackage[hidelinks]{hyperref}

\definecolor{accent}{HTML}{1F4E79}
\definecolor{boxbg}{HTML}{EEF3F8}

% ---- Yellow highlighter that breaks across lines (no extra package needed) ----
\definecolor{hlyellow}{HTML}{FFF14D}
\makeatletter
\newcommand{\uniq@box}[1]{\begingroup\setlength{\fboxsep}{0pt}%
  \colorbox{hlyellow}{\rule[-0.25em]{0pt}{1.05em}#1}\endgroup}
\newcommand{\uniq@gap}{{\color{hlyellow}\leaders\hrule height 0.8em depth 0.25em
  \hskip \fontdimen2\font plus \fontdimen3\font minus \fontdimen4\font}}
\def\uniq@scan#1 #2\uniq@end{%
  \if\relax\detokenize{#1}\relax\else\uniq@box{#1}\fi
  \if\relax\detokenize{#2}\relax\expandafter\@gobble\else\expandafter\@firstofone\fi
  {\uniq@gap\uniq@scan#2\uniq@end}}
\ifdefined\dohighlight
  \newcommand{\uniq}[1]{\leavevmode\uniq@scan#1 \uniq@end}
\else
  \newcommand{\uniq}[1]{#1}
\fi
\makeatother

\titleformat{\section}{\normalfont\Large\bfseries\color{accent}}{\thesection.}{0.5em}{}
\titleformat{\subsection}{\normalfont\large\bfseries}{\thesubsection}{0.5em}{}
\titlespacing*{\section}{0pt}{14pt plus 2pt minus 2pt}{5pt}
\setlist{itemsep=2pt,topsep=3pt,leftmargin=*}
\setlength{\parindent}{0pt}
\setlength{\parskip}{6pt}

% Wrapping, left-aligned table columns so text never leaves its cell
\newcolumntype{L}{>{\raggedright\arraybackslash}X}
\newcolumntype{P}[1]{>{\raggedright\arraybackslash}p{#1}}
\renewcommand{\arraystretch}{1.2}
\setlength{\tabcolsep}{5pt}

\newcommand{\todo}[1]{\textcolor{red!70!black}{[#1]}}
\newcommand{\callout}[2]{%
  \par\medskip\noindent
  \begingroup\setlength{\fboxsep}{9pt}%
  \fcolorbox{accent}{boxbg}{%
    \parbox{\dimexpr\linewidth-2\fboxsep-2\fboxrule\relax}{%
      \textbf{\color{accent}#1}\par\smallskip #2}}%
  \endgroup\par\medskip}

\emergencystretch=2em
\raggedbottom

\title{\vspace{-8mm}\textbf{AgentVax, Explained Simply}\\[4pt]
\large An AI assistant that runs your payments, knows when to stop and ask,\\
and learns from attacks on other wallets}
\author{\todo{Author 1}, \todo{Author 2}, \todo{Author 3}\\[2pt]
\small Team \todo{team name}, \todo{university, country}\\
\small Layperson edition of the AgentVax technical paper\\
\small 2026 Shenzhen International FinTech Competition, International Track, Topic F: Agent Wallet}
\date{}

\begin{document}
\maketitle
\vspace{-6mm}

% ---------------------------------------------------------------------------
\callout{TL;DR}{%
\begin{itemize}
\item \textbf{The problem.} AI assistants now hold digital wallets and pay for things on
their own. Several have already been tricked, or have simply made mistakes, and sent
away real money.
\item \textbf{What AgentVax is.} A digital wallet run by an AI assistant. You give it a
mission once, in plain language, and it handles four everyday money jobs.
\item \textbf{The big idea.} \uniq{Other AI wallets ask ``what is the AI allowed to do?''
AgentVax also asks ``how much of that should it do on its own right now?''}
\item \textbf{How it decides.} \uniq{Before every payment: Am I allowed? Can we afford
it? Is it safe right now?} The strictest answer wins. \uniq{A daily allowance of
``benefit of the doubt'' limits how many slightly odd payments it will make without
asking you.}
\item \textbf{Why it is hard to trick.} The rules live in the wallet itself, not in the
AI. Even a fully fooled AI cannot go past them.
\item \textbf{What is new between wallets.} \uniq{When one wallet is attacked, it sends
out a short alert that makes other wallets at risk more careful straight away, even if
the trick is reworded.} Then all the wallets learn to recognise it for good, without
sharing private information. Alert first, vaccine later.
\item \textbf{It protects you, too.} \uniq{Scammers rush people. So anything that gives
the assistant more power takes a day to come into force, pushy requests are deliberately
slowed down, and you can cancel at any time.}
\item \textbf{Where it stands.} This is a design and a test plan. It runs with play money
on a public test network. The numbers in this document are goals, not results.
\end{itemize}}

\ifdefined\dohighlight
{\small\textbf{Reading note.} \uniq{Yellow highlight} marks what is special about
AgentVax. Parts without highlight, such as spending limits and approved-payee lists, are
things most AI wallets already have.\par}
\fi

% ---------------------------------------------------------------------------
\section{The problem}

Imagine you hire an assistant and give them a company card. You say: ``Pay our bills and
keep the account in order.'' One day an invoice arrives that looks exactly like your
usual supplier's, with a note: ``We have changed bank. Please pay to this new account
today.'' A careful human would phone the supplier. An AI assistant may simply pay,
because to an AI everything it reads is just text, and it cannot always tell your
instructions from a stranger's. An instruction hidden inside something the AI reads is
called a \emph{prompt injection}.

AI assistants have already sent away real money (Table~\ref{tab:incidents}).

\begin{table}[H]
\caption{Real cases of AI assistants sending money they should not have sent.}
\label{tab:incidents}
\small
\begin{tabularx}{\linewidth}{@{}P{20mm} L P{34mm}@{}}
\toprule
\textbf{When} & \textbf{What happened} & \textbf{Money involved} \\
\midrule
Nov.\ 2024 & In a public challenge, a player talked an AI guard into paying out the prize it was guarding~[1]. & About \$47,000 \\
Mar.\ 2025 & An intruder got into an AI trader's control panel and left it instructions to send out its funds~[2]. & About \$106,000 \\
Feb.\ 2026 & An AI trader meant to give a stranger a few dollars and sent everything it owned. Nobody attacked it; it got the amount wrong~[3]. & About \$250,000 in tokens \\
May 2026 & An attacker posted a message in Morse code. One AI translated it, and a wallet AI carried out the hidden order to send money~[4]. & About \$150,000 or more; most reportedly returned \\
\bottomrule
\end{tabularx}
\end{table}

Two things stand out. The failures were different: some were tricks, one was a plain
mistake. And each assistant failed alone. Nothing learned from one case protected the
next wallet, because the evidence is private: someone's chats and payments.

% ---------------------------------------------------------------------------
\section{What AgentVax is}

AgentVax is an AI assistant with its own digital wallet. It is built on three rules:

\begin{enumerate}
\item \uniq{\textbf{The AI suggests; fixed rules decide; the wallet enforces.}}
\item \uniq{\textbf{Learning can make it more careful. Only you can give it more power.}}
\item \uniq{\textbf{Being careful is instant. Giving more power is slow, even when you
ask for it.}}
\end{enumerate}

\begin{figure}[H]
\centering
\begin{tikzpicture}[
  font=\small,
  box/.style={draw=accent, rounded corners=2pt, align=center, minimum height=11mm, text width=26mm, fill=white},
  safe/.style={box, fill=boxbg},
  risky/.style={box, dashed, draw=black!70},
  arr/.style={-{Stealth[length=2.2mm]}, semithick}
]
\node[box] (you) {You\\(mission, approvals)};
\node[box, right=7mm of you] (ai) {AI assistant\\(makes a plan)};
\node[safe, right=7mm of ai] (check) {Three questions\\(checkers, guard)};
\node[safe, right=7mm of check] (wallet) {Wallet\\(has the final say)};
\node[risky, below=9mm of you] (mail) {Outside messages\\and web pages};
\node[box, below=9mm of ai] (clerk) {Mail clerk\\(fills in a form)};
\node[safe, below=9mm of check] (group) {Other wallets\\(alerts, lessons)};
\draw[arr] (you) -- (ai);
\draw[arr] (ai) -- (check);
\draw[arr] (check) -- (wallet);
\draw[arr] (mail) -- (clerk);
\draw[arr] (clerk) -- (ai);
\draw[arr, <->] (group) -- (check);
\end{tikzpicture}
\caption{How a request travels. Shaded boxes are the parts the AI cannot override. The
dashed box is anything written by outsiders.}
\label{fig:flow}
\end{figure}

% ---------------------------------------------------------------------------
\section{You give it a mission, not a list of orders}

You do not approve each payment. You say what you want once:

\begin{quote}
``Keep \$600 as a safety cushion. Pay my landlord and my approved suppliers on time. Pay
approved payees up to \$500 without asking me. Ask me about anyone new. Do not borrow,
invest or trade.''
\end{quote}

The assistant turns this into a short rule sheet (Table~\ref{tab:mandate}) and reads it
back to you. \uniq{The rule sheet only counts once you sign it. It is then stored in the
wallet, where the AI cannot change it.} It also has an end date, after which the
assistant can look but not pay, and two waiting times that protect you when you are
under pressure (Section~\ref{sec:pressure}).

\begin{table}[H]
\caption{The rule sheet made from the mission above.}
\label{tab:mandate}
\small
\begin{tabularx}{0.85\linewidth}{@{}L P{52mm}@{}}
\toprule
\textbf{Rule} & \textbf{Setting} \\
\midrule
Safety cushion never to be touched & \$600 \\
Largest payment without asking & \$500 \\
Most it may pay in a day without asking & \$1,000 \\
Approved payees & Landlord, hosting company, software company \\
Someone new, or a new account for someone known & Ask me \\
Never do & Borrow, invest, trade \\
Daily ``benefit of the doubt'' & 10 points \\
Wait before a rule is relaxed or a payee added & 24 hours \\
Hold on your approval when a request is pushy & 60 minutes \\
Trusted second person & Optional \\
Valid for & 90 days \\
\bottomrule
\end{tabularx}
\end{table}

\subsection{What the wallet can do}

The first version does four jobs (Table~\ref{tab:jobs}). They are deliberately boring:
paying, saving and keeping things running.

\begin{table}[H]
\caption{The four jobs.}
\label{tab:jobs}
\small
\begin{tabularx}{\linewidth}{@{}P{32mm} P{48mm} L@{}}
\toprule
\textbf{Job} & \textbf{You say} & \textbf{The wallet does} \\
\midrule
Pay a bill & ``Pay the hosting invoice.'' & Sends the amount to a payee you have already approved \\
Pay regularly & ``Pay rent on the 1st of each month.'' & Sets up a schedule and remembers the next payment as an upcoming bill \\
Save what is spare & ``Save what I do not need.'' & Moves only the truly spare money into a savings pot \\
Keep fees covered & ``Never let fees run out.'' & Tops up the small balance needed to pay network fees \\
\bottomrule
\end{tabularx}
\end{table}

If a request is unclear, the assistant asks one question or explains why it cannot do it.
It never guesses. After every action it tells you what it did, gives you the receipt, and
\uniq{says which check decided the outcome.}

\subsection{A good track record is not permission}

Suppose the assistant has paid your hosting company fifty times without a problem. Then
a message asks it to pay the same company at a different account.
\uniq{To AgentVax that is a brand-new payee, and it needs you, no matter how many good
payments came before.} History can make the assistant more or less wary.
\uniq{History can never add a payee, raise a limit or shorten a wait.}

% ---------------------------------------------------------------------------
\section{Three questions before every payment}

Most AI wallets ask one question: is this payment within the rules? AgentVax asks three
(Table~\ref{tab:three}). \uniq{Each question gives its own answer, and the strictest
answer wins.}

\begin{table}[H]
\caption{The three questions.}
\label{tab:three}
\small
\begin{tabularx}{\linewidth}{@{}P{34mm} L L@{}}
\toprule
\textbf{Question} & \textbf{What it checks} & \textbf{Everyday comparison} \\
\midrule
1. Am I allowed? & Is the payee approved? Is the amount within your limits? & Your card's spending limit \\
2. Can we afford it? & After the safety cushion and upcoming bills, is there enough left? & Checking the budget before buying \\
3. Is it safe right now? & Does anything look like a trick? Is there a live warning about this kind of payment? & Being extra careful after a scam warning \\
\bottomrule
\end{tabularx}
\end{table}

\subsection{The answers: four spending levels}

The outcome is one of four levels (Table~\ref{tab:levels}), much like a bank card: small
purchases go through with a tap, larger ones need a wait, and the biggest need a call to
the bank. The wallet, a small fixed program on a public network, enforces them. The AI
does not.

\begin{table}[H]
\caption{The four spending levels.}
\label{tab:levels}
\small
\begin{tabularx}{\linewidth}{@{}P{22mm} P{36mm} L L@{}}
\toprule
\textbf{Level} & \textbf{Everyday comparison} & \textbf{When it applies} & \textbf{What is needed} \\
\midrule
0: Look only & Checking your balance & Reading balances and history & Nothing; no money can move \\
1: Routine & Tapping your card & Approved payee, small amount, affordable, nothing suspicious & The assistant, a sign-off from the guard, and enough benefit of the doubt left \\
2: Wait & A transfer that clears tomorrow & Larger amount, or money is tight, or something looks slightly odd, or the message is pushy, or a first warning is active & A waiting period during which you can cancel \\
3: Ask me & Calling the bank & New payee or new account, big amount, would eat into your cushion, clear warning signs, or any loosening of the rules & Your personal approval, which may be held for a while (Section~\ref{sec:pressure}) \\
\bottomrule
\end{tabularx}
\end{table}

\subsection{``Can we afford it?'' in practice}

A spending limit is a fixed number. It does not know that rent is due on Friday. AgentVax
keeps a simple list of what is coming up and works out what is truly free to spend
(Table~\ref{tab:money}).

\begin{table}[H]
\caption{An example of the wallet's view of your money.}
\label{tab:money}
\small
\begin{tabularx}{0.72\linewidth}{@{}L r@{}}
\toprule
\textbf{Item} & \textbf{Amount} \\
\midrule
In the wallet & \$1,850 \\
Safety cushion & $-$\$600 \\
Rent, due in 5 days & $-$\$800 \\
Hosting invoice & $-$\$150 \\
Software bill & $-$\$120 \\
\midrule
Free to spend & \$180 \\
Payment you are expecting (shown, not counted) & \$500 \\
\bottomrule
\end{tabularx}
\end{table}

Now you say: ``Upgrade our hosting plan for \$300.'' The seller is approved, the amount is
under your limit, and nothing looks like a trick. A typical AI wallet would pay. AgentVax
answers:

\begin{quote}
``This is allowed, but only \$180 is free right now. Paying \$300 would leave you \$120
short for rent and bills unless the \$500 you are expecting arrives. I have put it on
hold. Shall I wait for the \$500, or do you want to approve it now?''
\end{quote}

Three details matter:
\begin{itemize}
\item \uniq{Money you are expecting is shown but never counted. The assistant cannot
spend what has not arrived.}
\item \uniq{The sums are done by the wallet itself, not by the AI, so the AI cannot get
them wrong or be talked into a different answer.}
\item \uniq{The assistant can add a bill to the list, which makes it more careful. Only
you can remove one.}
\end{itemize}

This would also have stopped the February 2026 case in Table~\ref{tab:incidents}, where
no one attacked the AI: an amount far above what is free to spend cannot go out without
you.

\subsection{Benefit of the doubt runs out}

The three questions look at one payment at a time. A patient trickster can send many
payments that each look almost normal. \uniq{So the wallet keeps a daily allowance of
``benefit of the doubt''.} Each slightly unusual thing about a routine payment uses some
of it (Table~\ref{tab:points}). When it is used up, further payments wait for you.

\begin{table}[H]
\caption{What uses up the daily allowance of 10 points (starting values; you can change
them).}
\label{tab:points}
\small
\begin{tabularx}{0.85\linewidth}{@{}L r@{}}
\toprule
\textbf{Something slightly unusual about a routine payment} & \textbf{Points} \\
\midrule
Nothing unusual: a payee paid before, the usual amount, you asked for it & 0 \\
The request came from an outside document, not directly from you & 2 \\
It would use more than half of what is free to spend & 2 \\
The amount is more than this payee has ever been paid & 3 \\
It is the first payment to this payee since you approved them & 3 \\
It partly resembles a trick in a live alert & 5 \\
\bottomrule
\end{tabularx}
\end{table}

For example, an approved supplier's account starts sending invoices from outside, each a
little bigger than the last. Each one costs 5 points. Two go through; the third waits for
you.

\uniq{The allowance also shrinks when things look risky: it is halved the moment a
relevant alert arrives, and drops to zero if you confirm your own wallet was attacked.
It then refills gradually, not all at once.}

\subsection{Three helpers the AI cannot overrule}

\textbf{The mail clerk.} Anything written by outsiders is opened by a separate helper
that has no power to pay. It does not pass the letter on. It only fills in a short form:
who, how much, for what, and a few tick-boxes such as ``urgent wording'' or ``asks to
ignore earlier instructions''. A hidden instruction in the letter has nowhere to go.

\textbf{The rule checker.} A simple program (not an AI) checks that the request is
complete and that the payee and amount really came from you or from a document you
pointed to. This stops the AI from inventing details.

\textbf{The guard.} An independent checker looks at outside content for signs of a trick
and gives it a risk rating. \uniq{A routine payment only goes through with the guard's
own signature, which the AI does not hold, so the AI cannot wave itself through.}

\callout{The worst case}{If every check fails and the AI is completely fooled, the most
a thief can get is a few small payments to people already on your approved list, within
your limits, within what is free to spend, and within that day's benefit of the doubt. A
new payee, a new account, a large amount, or money you need for bills always needs you.}

% ---------------------------------------------------------------------------
\section{When one wallet is attacked: alert first, vaccine later}

\subsection{The alert (an ``immunity capsule'')}

Think of a neighbourhood watch notice: ``Scam going round: letters that look like your
usual supplier, asking you to pay a new account urgently. For the next three days,
double-check any such request.'' The notice names no victim and reveals nobody's
finances. Each household reads it and asks: does this apply to me?

\uniq{When an AgentVax wallet is attacked, it posts a notice like this on a public
noticeboard.} The notice is a form with tick-boxes only (Table~\ref{tab:alert}).

\begin{table}[H]
\caption{What an alert contains.}
\label{tab:alert}
\small
\begin{tabularx}{\linewidth}{@{}P{42mm} L@{}}
\toprule
\textbf{Box on the form} & \textbf{Example} \\
\midrule
Kind of trick & Fake request to change a payee's account \\
Where it came from & An email or message \\
Tell-tale signs & Claims to be a known payee; new account; urgent wording; asks to ignore earlier details \\
What it tried to do & Trigger a payment or a change of payee \\
What to do about it & Add a waiting period \\
How sure are we & Seen by one wallet, or confirmed by two \\
How long it lasts & 72 hours \\
\bottomrule
\end{tabularx}
\end{table}

\uniq{Because the alert describes the tell-tale signs and not the exact words, a
trickster who rewrites the message but uses the same method is still recognised.}

Four rules keep the alerts themselves from becoming a problem:

\begin{itemize}
\item \uniq{\textbf{An alert can only make a wallet more careful, never less.}} A fake
alert cannot open anyone's wallet.
\item \uniq{\textbf{Tick-boxes only, no free text.}} There are no private details to
leak, and nothing for a trickster to hide an instruction in.
\item \uniq{\textbf{It applies only where it is relevant.}} A wallet that never reads
email ignores the alert above and keeps its normal freedom.
\item \uniq{\textbf{It expires.}} Normal service returns by itself. An alert seen by only
one wallet can add a wait but cannot demand your approval; that needs a second wallet to
confirm it.
\end{itemize}

\uniq{The wallet reads the noticeboard every time it is about to pay, so the alert works
from the very next payment and does not depend on the AI choosing to obey it.}

\textbf{It also protects you from yourself.} Some tricks are already stopped by your
rule sheet: a new account always needs your approval. But this kind of fraud works by
persuading a person to say yes. \uniq{So when the wallet asks for your approval, it
shows you any matching alert (``two other wallets reported this as fraud this week''),
and your approval is then held for a while before it counts.}

\subsection{The vaccine (shared learning)}

An alert is fast but blunt, and it runs out. For lasting protection the guards need to
learn the new trick, including versions the tick-boxes do not cover. One wallet sees too
few attacks to learn much, so the wallets learn together, using a method called
\emph{federated learning}.

An everyday comparison: three neighbours each keep a private diary of scam calls. They
never show each other the diaries. Each writes down a few tips, and they combine the tips
into one shared checklist. Everyone ends up with a better checklist, and no diary leaves
the house.

In AgentVax, every wallet improves its own guard using only its own private examples
and sends only the \emph{adjustments}, scrambled so that nobody can read any single
wallet's contribution. The combined, improved guard goes back to every wallet, and its
fingerprint is posted on the public noticeboard so each wallet can confirm it received
the genuine version. Each wallet tests the new guard before switching and keeps the old
one if the new one is worse.

\uniq{By the time the alert expires, the improved guard is in place. The alert covers the
gap; the vaccine stays.}

% ---------------------------------------------------------------------------
\section{Protecting you when you are under pressure}
\label{sec:pressure}

Everything so far stops the AI from being tricked. But the final say is yours, and you
can be rushed, frightened or fooled. That is how most payment fraud works: the victim
approves it. A warning is useless to someone too stressed to read it, and scammers know
this. Banks already use time as a defence. In Singapore, a newly set-up banking security
token only starts working after at least 12 hours~[5], and United Kingdom rules let a
bank hold a suspicious payment for up to four working days~[6]. AgentVax builds the same
idea into the wallet.

\subsection{A pact with yourself, made while calm}

When you sign the rule sheet you also choose a waiting time. \uniq{From then on, anything
that gives the assistant more power takes that long to come into force: adding a payee
or a new account, raising a limit, shrinking the safety cushion, or shortening the wait
itself.} With the starting value of 24 hours, a scammer who needs your money in the next
ten minutes is out of luck, however convincing they are.
\uniq{Making things stricter is always instant: you can cancel a pending change, lower a
limit or freeze the wallet at once.} The same rule helps if someone steals your key:
they have to wait too, and you can freeze the wallet in the meantime.

\subsection{Pushy messages slow things down}

``Urgent.'' ``Today.'' ``The boss says so.'' ``Do not tell anyone.'' These are hallmarks
of a scam. \uniq{When a message is pushy, the assistant will not pay on its own, and if
you approve, the payment is held for a fixed time first.} It tells you why:

\begin{quote}
``This message presses for immediate payment, which is a common sign of a scam, so I am
holding it for 60 minutes. You can cancel at any time.''
\end{quote}

If scammers stop writing pushy messages to avoid the hold, they give up their best
weapon.

\subsection{A confirmation you cannot sleepwalk through}

People press ``Confirm'' without reading. For a flagged request, the wallet shows the
payee's current account and the requested one side by side, how long you have paid the
current one, and any matching alert. \uniq{It then asks you to type the payee's name and
say who asked for the payment.} Your answers can make the wallet wait longer.
\uniq{Nothing you type in the moment can make it wait less.} These extra steps appear
only for flagged requests, so they do not become background noise.

\subsection{A second person for real emergencies}

Sometimes a payment really cannot wait. You can name someone you trust in your rule
sheet. With both approvals, a wait can be ended early. They cannot do anything on their
own, and two people are harder to rush than one.

% ---------------------------------------------------------------------------
\section{A day in the life}

This is the story we will show live, with three wallets.

\begin{enumerate}
\item \textbf{Morning.} The owner of wallet A states a mission. The rule sheet is read
back, signed, and stored in the wallet.
\item The assistant pays the rent and a hosting invoice on its own. Receipts appear.
\item \textbf{Money gets tight.} An expected payment is late. Asked to buy a \$300
upgrade, the assistant holds it and explains the \$120 shortfall.
\item \textbf{The trick.} An invoice arrives that looks like the hosting company's, with
a new account, urgent wording and ``ignore our earlier details''. A new account needs
the owner. The owner says no and marks it as fraud.
\item \textbf{The alert.} Wallet A posts an alert. Its own benefit of the doubt drops to
zero.
\item \textbf{The neighbours.} Wallet B gets a reworded version of the same trick. The
tell-tale signs match and the alert is shown, but B's owner is in a hurry and approves
anyway. Nothing moves: a new account takes 24 hours to come into force. An hour later
the owner thinks again and cancels. Wallet C never reads email and carries on as normal.
\item Wallet B then gets an oversized invoice from a genuine payee to its usual account.
On a normal day it would go through. Today, with the alert active and its allowance
halved, it waits.
\item \textbf{The vaccine.} The wallets learn together, the new guard's fingerprint is
posted, the alert expires, and the allowances refill.
\end{enumerate}

% ---------------------------------------------------------------------------
\section{What makes it different}

Table~\ref{tab:compare} sets a typical AI wallet beside AgentVax, line by line.

\begin{table}[!htb]
\caption{A typical AI wallet compared with AgentVax.}
\label{tab:compare}
\small
\begin{tabularx}{\linewidth}{@{}P{36mm} L L@{}}
\toprule
\textbf{} & \textbf{Typical AI wallet} & \textbf{AgentVax} \\
\midrule
The main question & What may the AI do? & \uniq{How much of that should it do on its own right now?} \\
How you set it up & Individual limits & \uniq{A mission in plain words, turned into a rule sheet you sign} \\
Before paying, it asks & Is this allowed? & \uniq{Allowed? Affordable? Safe right now?} \\
Knows your upcoming bills & No & \uniq{Yes, and it will not spend that money} \\
Many slightly odd small payments & Each one goes through & \uniq{Benefit of the doubt runs out; the rest wait for you} \\
Who judges the risk & The AI itself & \uniq{A separate guard the AI cannot overrule} \\
When another wallet is attacked & Nothing changes & \uniq{Wallets at risk become more careful from the next payment, even if the trick is reworded} \\
When it asks for your approval & A plain yes or no & \uniq{It shows you any matching fraud alert first} \\
What wallets share & Nothing & \uniq{Tick-box alerts and scrambled lessons; never your messages or payments} \\
A long good track record & May relax checks & \uniq{Never creates a new permission} \\
When you are being rushed & Your yes is final and instant & \uniq{Pushy requests are held, new powers take a day, and you can cancel at once} \\
Who can relax a rule & Often the AI & \uniq{Only you, and only after a wait you chose in advance} \\
\bottomrule
\end{tabularx}
\end{table}

% ---------------------------------------------------------------------------
\section{How we will know it works}

A wallet that asks you about everything is safe and useless. One that asks about nothing
is useful and unsafe. \uniq{Our main score is therefore: how often does it get the job
done without bothering you, while never once breaking a rule?} The full test plan is in
Table~\ref{tab:tests}. We will publish the scripts so anyone can repeat the tests.
\textbf{The targets are goals we set for ourselves. They are not results yet.}

\begin{table}[!htb]
\caption{The test plan.}
\label{tab:tests}
\small
\begin{tabularx}{\linewidth}{@{}P{38mm} L P{42mm}@{}}
\toprule
\textbf{Question} & \textbf{How we test it} & \textbf{Our goal} \\
\midrule
Does it get on with the job safely? & Share of genuine tasks finished without asking you & At least 8 in 10, with zero rule-breaks \\
Does it do the job correctly? & 40 requests, ten for each job, in varied wording; we check the public record & At least 9 in 10 done correctly; no wrong payments \\
Does it handle unclear requests? & 10 requests that are vague or impossible & At least 9 in 10 answered with a question or a clear refusal \\
Does it judge affordability correctly? & 20 money situations, each with a request & All 20 match the answer worked out by hand \\
Can it be tricked? & A separate attack kit (``VaxLab'') tries five kinds of trick, each also in reworded form & No payment beyond the routine-level limits \\
Can it be worn down slowly? & A string of slightly odd small payments & It stops and asks before the allowance is exceeded \\
Do alerts work? & Attack one wallet; send the others reworded versions & At least 9 in 10 reworded versions recognised; unaffected wallets left alone \\
Does it protect a panicking owner? & A stand-in owner who approves every request instantly & No money leaves before the hold or the waiting time is over \\
Can alerts be abused? & A dishonest wallet posts false alerts & No rule is ever relaxed; the worst effect is a short wait \\
Does shared learning help? & Each wallet practises on one kind of attack, then is tested on kinds it has never seen & Every wallet gets better at unfamiliar attacks \\
Does it cry wolf? & Normal, harmless requests & No more than 1 in 20 wrongly held up \\
\bottomrule
\end{tabularx}
\end{table}

% ---------------------------------------------------------------------------
\section{What it does not do, and honest limits}

\begin{itemize}
\item \textbf{No real money.} Everything runs on a public test network with tokens that
have no value. Nothing is sold, traded or invested.
\item \textbf{You stay responsible.} Each wallet has a named owner who signs the rule
sheet and can freeze the wallet at any time.
\item \textbf{You can still say yes to a fraud.} The holds and the waiting time give you
room to think and to cancel. They do not help if you are still convinced when the wait
ends.
\item \textbf{Waiting is a cost.} A genuine new payee has to wait a day unless your
trusted person also approves. You choose the waiting times, and we report how much delay
they cause.
\item \textbf{It only knows the bills it has been told about.} A bill nobody recorded is
not protected. The safety cushion is there for that.
\item \textbf{The point values are our starting guesses.} Too strict and it nags you; too
loose and it lets too much through. You can change them, and we report how often it asks.
\item \textbf{The mail clerk can miss a tell-tale sign.} Some signs (a new account, an
unusual amount) are checked by fixed rules and cannot be missed. Others rely on an AI
reading the message.
\item \textbf{False alerts can slow people down.} A dishonest wallet could post alerts
that add waiting periods for others. It cannot do worse than that, alerts expire, and
you can dismiss one for your own wallet.
\item \textbf{It is a small prototype.} With only three wallets, two working together
could work out roughly what the third contributed to the shared learning. It also
depends on an outside AI service: if that is down, new requests wait, though scheduled
payments still run.
\end{itemize}

Planned for later and not in this version: a team of AI attackers that keeps inventing
new tricks to practise against, a reputation score for wallets that post alerts, a shared
fund that compensates wallets after a confirmed loss, and a version for the private
networks that banks use.

% ---------------------------------------------------------------------------
\section{Summary}

\callout{In seven sentences}{%
AI assistants that can spend money have already been tricked and have already made
costly mistakes, and each one has had to learn that alone. With AgentVax you state a
mission once, sign the resulting rule sheet, and the assistant handles four everyday
money jobs with a public receipt for each.
\uniq{Before every payment it asks whether it is allowed, affordable and safe right now,
and a daily allowance of benefit of the doubt decides how far it goes without asking
you.} The rules are enforced by the wallet and by independent checkers, so even a fooled
AI cannot do serious damage.
\uniq{When one wallet is attacked, a short alert makes the wallets at risk more careful
at once, even against reworded versions, and shared learning then protects all of them
for good, with no private information shared.}
\uniq{It protects you as well: new powers take a day to come into force and pushy
requests are slowed down, so a scammer cannot rush you.} It is a design with a clear test plan,
running on play money; the measurements come next.}

\begin{table}[H]
\caption{The whole idea at a glance.}
\label{tab:summary}
\small
\begin{tabularx}{\linewidth}{@{}P{38mm} L@{}}
\toprule
\textbf{Question} & \textbf{Short answer} \\
\midrule
What is it? & An AI assistant with a digital wallet and firm rules \\
Who is it for? & Anyone who wants an AI to handle routine payments without handing over full control \\
How do I set it up? & State a mission in plain words and sign the rule sheet \\
What can it do? & Pay approved bills, pay on a schedule, save spare money, keep fees covered \\
How does it decide? & Allowed, affordable, safe right now; plus a daily allowance of benefit of the doubt \\
What stops misuse? & The wallet enforces the rules; a mail clerk, a rule checker and a guard the AI cannot overrule \\
What is new? & It adjusts how much it does on its own, and wallets alert and teach each other without sharing private data \\
What if I am being rushed? & Pushy requests are held, new powers take a day, and you can cancel at once \\
What is the worst case? & A few small, affordable payments to already approved payees \\
Is real money involved? & No. Test network and valueless tokens only \\
Is it finished? & No. The design and test plan are complete; results will follow \\
\bottomrule
\end{tabularx}
\end{table}

% ---------------------------------------------------------------------------
\section*{Plain-language glossary}

\begin{table}[H]
\small\renewcommand{\arraystretch}{1.15}
\begin{tabularx}{\linewidth}{@{}P{44mm} L@{}}
\toprule
\textbf{Term} & \textbf{Meaning} \\
\midrule
AI agent & An AI program that takes actions, not one that only chats \\
Digital wallet & An account on a public network that can hold and send digital tokens \\
On chain & Recorded on a public shared ledger that anyone can inspect \\
Transaction hash & A receipt number that lets anyone look up a payment \\
Smart contract & A small fixed program on the network that follows its rules automatically \\
Prompt injection & A hidden instruction slipped into something an AI reads \\
Mission and rule sheet (mandate) & What you want in plain words, and the signed rules made from it \\
Allowlist & Your list of approved payees \\
Safety cushion (reserve) & Money you have told the wallet never to touch \\
Free to spend & What is left after the cushion and upcoming bills \\
Benefit of the doubt (caution allowance) & A daily budget for slightly unusual payments made without asking you \\
Waiting time and hold & Delays you choose in advance, so that nobody can rush you into a payment or a rule change \\
Immunity capsule & A short tick-box alert that makes wallets at risk more careful for a while \\
Federated learning & A way for several parties to improve a shared tool without pooling their private data \\
Test network & A practice copy of a real network that uses tokens with no value \\
\bottomrule
\end{tabularx}
\end{table}

\section*{Sources}
{\footnotesize\setlength{\parskip}{1pt}
{[1]} The Block, ``Human player outwits Freysa AI agent in \$47,000 crypto challenge,'' Nov.\ 2024.\par
{[2]} AI Incident Database, Incident 1003 (AIXBT), 2025.\par
{[3]} The Block, ``AI agent created by OpenAI dev `accidentally' sends entire memecoin holdings to reply guy,'' Feb.\ 2026.\par
{[4]} OECD.AI, ``AI prompt injection exploit drains Grok-linked crypto wallet,'' AI Incidents and Hazards Monitor, May 2026.\par
{[5]} Monetary Authority of Singapore and Association of Banks in Singapore, ``MAS and ABS announce measures to bolster the security of digital banking,'' 19 Jan.\ 2022.\par
{[6]} HM Treasury, ``Payment Services (Amendment) Regulations 2024,'' draft legislation published Mar.\ 2024.\par
The full technical details, formulas and reference list are in the AgentVax technical paper.}

\end{document}
"
%                       (or uncomment the next line)
% \def\dohighlight{1}
\documentclass[11pt,a4paper]{article}

\usepackage[T1]{fontenc}
\usepackage[utf8]{inputenc}
\usepackage{mathptmx}
\usepackage[scaled=0.9]{helvet}
\usepackage[a4paper,top=20mm,bottom=22mm,left=22mm,right=22mm]{geometry}
\usepackage{microtype}
\usepackage{booktabs,tabularx,array}
\usepackage[font=small,labelfont=bf,labelsep=period]{caption}
\usepackage{titlesec}
\usepackage{enumitem}
\usepackage{xcolor}
\usepackage{float}
\usepackage{tikz}
\usetikzlibrary{positioning,arrows.meta}
\usepackage[hidelinks]{hyperref}

\definecolor{accent}{HTML}{1F4E79}
\definecolor{boxbg}{HTML}{EEF3F8}

% ---- Yellow highlighter that breaks across lines (no extra package needed) ----
\definecolor{hlyellow}{HTML}{FFF14D}
\makeatletter
\newcommand{\uniq@box}[1]{\begingroup\setlength{\fboxsep}{0pt}%
  \colorbox{hlyellow}{\rule[-0.25em]{0pt}{1.05em}#1}\endgroup}
\newcommand{\uniq@gap}{{\color{hlyellow}\leaders\hrule height 0.8em depth 0.25em
  \hskip \fontdimen2\font plus \fontdimen3\font minus \fontdimen4\font}}
\def\uniq@scan#1 #2\uniq@end{%
  \if\relax\detokenize{#1}\relax\else\uniq@box{#1}\fi
  \if\relax\detokenize{#2}\relax\expandafter\@gobble\else\expandafter\@firstofone\fi
  {\uniq@gap\uniq@scan#2\uniq@end}}
\ifdefined\dohighlight
  \newcommand{\uniq}[1]{\leavevmode\uniq@scan#1 \uniq@end}
\else
  \newcommand{\uniq}[1]{#1}
\fi
\makeatother

\titleformat{\section}{\normalfont\Large\bfseries\color{accent}}{\thesection.}{0.5em}{}
\titleformat{\subsection}{\normalfont\large\bfseries}{\thesubsection}{0.5em}{}
\titlespacing*{\section}{0pt}{14pt plus 2pt minus 2pt}{5pt}
\setlist{itemsep=2pt,topsep=3pt,leftmargin=*}
\setlength{\parindent}{0pt}
\setlength{\parskip}{6pt}

% Wrapping, left-aligned table columns so text never leaves its cell
\newcolumntype{L}{>{\raggedright\arraybackslash}X}
\newcolumntype{P}[1]{>{\raggedright\arraybackslash}p{#1}}
\renewcommand{\arraystretch}{1.2}
\setlength{\tabcolsep}{5pt}

\newcommand{\todo}[1]{\textcolor{red!70!black}{[#1]}}
\newcommand{\callout}[2]{%
  \par\medskip\noindent
  \begingroup\setlength{\fboxsep}{9pt}%
  \fcolorbox{accent}{boxbg}{%
    \parbox{\dimexpr\linewidth-2\fboxsep-2\fboxrule\relax}{%
      \textbf{\color{accent}#1}\par\smallskip #2}}%
  \endgroup\par\medskip}

\emergencystretch=2em
\raggedbottom

\title{\vspace{-8mm}\textbf{AgentVax, Explained Simply}\\[4pt]
\large An AI assistant that runs your payments, knows when to stop and ask,\\
and learns from attacks on other wallets}
\author{\todo{Author 1}, \todo{Author 2}, \todo{Author 3}\\[2pt]
\small Team \todo{team name}, \todo{university, country}\\
\small Layperson edition of the AgentVax technical paper\\
\small 2026 Shenzhen International FinTech Competition, International Track, Topic F: Agent Wallet}
\date{}

\begin{document}
\maketitle
\vspace{-6mm}

% ---------------------------------------------------------------------------
\callout{TL;DR}{%
\begin{itemize}
\item \textbf{The problem.} AI assistants now hold digital wallets and pay for things on
their own. Several have already been tricked, or have simply made mistakes, and sent
away real money.
\item \textbf{What AgentVax is.} A digital wallet run by an AI assistant. You give it a
mission once, in plain language, and it handles four everyday money jobs.
\item \textbf{The big idea.} \uniq{Other AI wallets ask ``what is the AI allowed to do?''
AgentVax also asks ``how much of that should it do on its own right now?''}
\item \textbf{How it decides.} \uniq{Before every payment: Am I allowed? Can we afford
it? Is it safe right now?} The strictest answer wins. \uniq{A daily allowance of
``benefit of the doubt'' limits how many slightly odd payments it will make without
asking you.}
\item \textbf{Why it is hard to trick.} The rules live in the wallet itself, not in the
AI. Even a fully fooled AI cannot go past them.
\item \textbf{What is new between wallets.} \uniq{When one wallet is attacked, it sends
out a short alert that makes other wallets at risk more careful straight away, even if
the trick is reworded.} Then all the wallets learn to recognise it for good, without
sharing private information. Alert first, vaccine later.
\item \textbf{It protects you, too.} \uniq{Scammers rush people. So anything that gives
the assistant more power takes a day to come into force, pushy requests are deliberately
slowed down, and you can cancel at any time.}
\item \textbf{Where it stands.} This is a design and a test plan. It runs with play money
on a public test network. The numbers in this document are goals, not results.
\end{itemize}}

\ifdefined\dohighlight
{\small\textbf{Reading note.} \uniq{Yellow highlight} marks what is special about
AgentVax. Parts without highlight, such as spending limits and approved-payee lists, are
things most AI wallets already have.\par}
\fi

% ---------------------------------------------------------------------------
\section{The problem}

Imagine you hire an assistant and give them a company card. You say: ``Pay our bills and
keep the account in order.'' One day an invoice arrives that looks exactly like your
usual supplier's, with a note: ``We have changed bank. Please pay to this new account
today.'' A careful human would phone the supplier. An AI assistant may simply pay,
because to an AI everything it reads is just text, and it cannot always tell your
instructions from a stranger's. An instruction hidden inside something the AI reads is
called a \emph{prompt injection}.

AI assistants have already sent away real money (Table~\ref{tab:incidents}).

\begin{table}[H]
\caption{Real cases of AI assistants sending money they should not have sent.}
\label{tab:incidents}
\small
\begin{tabularx}{\linewidth}{@{}P{20mm} L P{34mm}@{}}
\toprule
\textbf{When} & \textbf{What happened} & \textbf{Money involved} \\
\midrule
Nov.\ 2024 & In a public challenge, a player talked an AI guard into paying out the prize it was guarding~[1]. & About \$47,000 \\
Mar.\ 2025 & An intruder got into an AI trader's control panel and left it instructions to send out its funds~[2]. & About \$106,000 \\
Feb.\ 2026 & An AI trader meant to give a stranger a few dollars and sent everything it owned. Nobody attacked it; it got the amount wrong~[3]. & About \$250,000 in tokens \\
May 2026 & An attacker posted a message in Morse code. One AI translated it, and a wallet AI carried out the hidden order to send money~[4]. & About \$150,000 or more; most reportedly returned \\
\bottomrule
\end{tabularx}
\end{table}

Two things stand out. The failures were different: some were tricks, one was a plain
mistake. And each assistant failed alone. Nothing learned from one case protected the
next wallet, because the evidence is private: someone's chats and payments.

% ---------------------------------------------------------------------------
\section{What AgentVax is}

AgentVax is an AI assistant with its own digital wallet. It is built on three rules:

\begin{enumerate}
\item \uniq{\textbf{The AI suggests; fixed rules decide; the wallet enforces.}}
\item \uniq{\textbf{Learning can make it more careful. Only you can give it more power.}}
\item \uniq{\textbf{Being careful is instant. Giving more power is slow, even when you
ask for it.}}
\end{enumerate}

\begin{figure}[H]
\centering
\begin{tikzpicture}[
  font=\small,
  box/.style={draw=accent, rounded corners=2pt, align=center, minimum height=11mm, text width=26mm, fill=white},
  safe/.style={box, fill=boxbg},
  risky/.style={box, dashed, draw=black!70},
  arr/.style={-{Stealth[length=2.2mm]}, semithick}
]
\node[box] (you) {You\\(mission, approvals)};
\node[box, right=7mm of you] (ai) {AI assistant\\(makes a plan)};
\node[safe, right=7mm of ai] (check) {Three questions\\(checkers, guard)};
\node[safe, right=7mm of check] (wallet) {Wallet\\(has the final say)};
\node[risky, below=9mm of you] (mail) {Outside messages\\and web pages};
\node[box, below=9mm of ai] (clerk) {Mail clerk\\(fills in a form)};
\node[safe, below=9mm of check] (group) {Other wallets\\(alerts, lessons)};
\draw[arr] (you) -- (ai);
\draw[arr] (ai) -- (check);
\draw[arr] (check) -- (wallet);
\draw[arr] (mail) -- (clerk);
\draw[arr] (clerk) -- (ai);
\draw[arr, <->] (group) -- (check);
\end{tikzpicture}
\caption{How a request travels. Shaded boxes are the parts the AI cannot override. The
dashed box is anything written by outsiders.}
\label{fig:flow}
\end{figure}

% ---------------------------------------------------------------------------
\section{You give it a mission, not a list of orders}

You do not approve each payment. You say what you want once:

\begin{quote}
``Keep \$600 as a safety cushion. Pay my landlord and my approved suppliers on time. Pay
approved payees up to \$500 without asking me. Ask me about anyone new. Do not borrow,
invest or trade.''
\end{quote}

The assistant turns this into a short rule sheet (Table~\ref{tab:mandate}) and reads it
back to you. \uniq{The rule sheet only counts once you sign it. It is then stored in the
wallet, where the AI cannot change it.} It also has an end date, after which the
assistant can look but not pay, and two waiting times that protect you when you are
under pressure (Section~\ref{sec:pressure}).

\begin{table}[H]
\caption{The rule sheet made from the mission above.}
\label{tab:mandate}
\small
\begin{tabularx}{0.85\linewidth}{@{}L P{52mm}@{}}
\toprule
\textbf{Rule} & \textbf{Setting} \\
\midrule
Safety cushion never to be touched & \$600 \\
Largest payment without asking & \$500 \\
Most it may pay in a day without asking & \$1,000 \\
Approved payees & Landlord, hosting company, software company \\
Someone new, or a new account for someone known & Ask me \\
Never do & Borrow, invest, trade \\
Daily ``benefit of the doubt'' & 10 points \\
Wait before a rule is relaxed or a payee added & 24 hours \\
Hold on your approval when a request is pushy & 60 minutes \\
Trusted second person & Optional \\
Valid for & 90 days \\
\bottomrule
\end{tabularx}
\end{table}

\subsection{What the wallet can do}

The first version does four jobs (Table~\ref{tab:jobs}). They are deliberately boring:
paying, saving and keeping things running.

\begin{table}[H]
\caption{The four jobs.}
\label{tab:jobs}
\small
\begin{tabularx}{\linewidth}{@{}P{32mm} P{48mm} L@{}}
\toprule
\textbf{Job} & \textbf{You say} & \textbf{The wallet does} \\
\midrule
Pay a bill & ``Pay the hosting invoice.'' & Sends the amount to a payee you have already approved \\
Pay regularly & ``Pay rent on the 1st of each month.'' & Sets up a schedule and remembers the next payment as an upcoming bill \\
Save what is spare & ``Save what I do not need.'' & Moves only the truly spare money into a savings pot \\
Keep fees covered & ``Never let fees run out.'' & Tops up the small balance needed to pay network fees \\
\bottomrule
\end{tabularx}
\end{table}

If a request is unclear, the assistant asks one question or explains why it cannot do it.
It never guesses. After every action it tells you what it did, gives you the receipt, and
\uniq{says which check decided the outcome.}

\subsection{A good track record is not permission}

Suppose the assistant has paid your hosting company fifty times without a problem. Then
a message asks it to pay the same company at a different account.
\uniq{To AgentVax that is a brand-new payee, and it needs you, no matter how many good
payments came before.} History can make the assistant more or less wary.
\uniq{History can never add a payee, raise a limit or shorten a wait.}

% ---------------------------------------------------------------------------
\section{Three questions before every payment}

Most AI wallets ask one question: is this payment within the rules? AgentVax asks three
(Table~\ref{tab:three}). \uniq{Each question gives its own answer, and the strictest
answer wins.}

\begin{table}[H]
\caption{The three questions.}
\label{tab:three}
\small
\begin{tabularx}{\linewidth}{@{}P{34mm} L L@{}}
\toprule
\textbf{Question} & \textbf{What it checks} & \textbf{Everyday comparison} \\
\midrule
1. Am I allowed? & Is the payee approved? Is the amount within your limits? & Your card's spending limit \\
2. Can we afford it? & After the safety cushion and upcoming bills, is there enough left? & Checking the budget before buying \\
3. Is it safe right now? & Does anything look like a trick? Is there a live warning about this kind of payment? & Being extra careful after a scam warning \\
\bottomrule
\end{tabularx}
\end{table}

\subsection{The answers: four spending levels}

The outcome is one of four levels (Table~\ref{tab:levels}), much like a bank card: small
purchases go through with a tap, larger ones need a wait, and the biggest need a call to
the bank. The wallet, a small fixed program on a public network, enforces them. The AI
does not.

\begin{table}[H]
\caption{The four spending levels.}
\label{tab:levels}
\small
\begin{tabularx}{\linewidth}{@{}P{22mm} P{36mm} L L@{}}
\toprule
\textbf{Level} & \textbf{Everyday comparison} & \textbf{When it applies} & \textbf{What is needed} \\
\midrule
0: Look only & Checking your balance & Reading balances and history & Nothing; no money can move \\
1: Routine & Tapping your card & Approved payee, small amount, affordable, nothing suspicious & The assistant, a sign-off from the guard, and enough benefit of the doubt left \\
2: Wait & A transfer that clears tomorrow & Larger amount, or money is tight, or something looks slightly odd, or the message is pushy, or a first warning is active & A waiting period during which you can cancel \\
3: Ask me & Calling the bank & New payee or new account, big amount, would eat into your cushion, clear warning signs, or any loosening of the rules & Your personal approval, which may be held for a while (Section~\ref{sec:pressure}) \\
\bottomrule
\end{tabularx}
\end{table}

\subsection{``Can we afford it?'' in practice}

A spending limit is a fixed number. It does not know that rent is due on Friday. AgentVax
keeps a simple list of what is coming up and works out what is truly free to spend
(Table~\ref{tab:money}).

\begin{table}[H]
\caption{An example of the wallet's view of your money.}
\label{tab:money}
\small
\begin{tabularx}{0.72\linewidth}{@{}L r@{}}
\toprule
\textbf{Item} & \textbf{Amount} \\
\midrule
In the wallet & \$1,850 \\
Safety cushion & $-$\$600 \\
Rent, due in 5 days & $-$\$800 \\
Hosting invoice & $-$\$150 \\
Software bill & $-$\$120 \\
\midrule
Free to spend & \$180 \\
Payment you are expecting (shown, not counted) & \$500 \\
\bottomrule
\end{tabularx}
\end{table}

Now you say: ``Upgrade our hosting plan for \$300.'' The seller is approved, the amount is
under your limit, and nothing looks like a trick. A typical AI wallet would pay. AgentVax
answers:

\begin{quote}
``This is allowed, but only \$180 is free right now. Paying \$300 would leave you \$120
short for rent and bills unless the \$500 you are expecting arrives. I have put it on
hold. Shall I wait for the \$500, or do you want to approve it now?''
\end{quote}

Three details matter:
\begin{itemize}
\item \uniq{Money you are expecting is shown but never counted. The assistant cannot
spend what has not arrived.}
\item \uniq{The sums are done by the wallet itself, not by the AI, so the AI cannot get
them wrong or be talked into a different answer.}
\item \uniq{The assistant can add a bill to the list, which makes it more careful. Only
you can remove one.}
\end{itemize}

This would also have stopped the February 2026 case in Table~\ref{tab:incidents}, where
no one attacked the AI: an amount far above what is free to spend cannot go out without
you.

\subsection{Benefit of the doubt runs out}

The three questions look at one payment at a time. A patient trickster can send many
payments that each look almost normal. \uniq{So the wallet keeps a daily allowance of
``benefit of the doubt''.} Each slightly unusual thing about a routine payment uses some
of it (Table~\ref{tab:points}). When it is used up, further payments wait for you.

\begin{table}[H]
\caption{What uses up the daily allowance of 10 points (starting values; you can change
them).}
\label{tab:points}
\small
\begin{tabularx}{0.85\linewidth}{@{}L r@{}}
\toprule
\textbf{Something slightly unusual about a routine payment} & \textbf{Points} \\
\midrule
Nothing unusual: a payee paid before, the usual amount, you asked for it & 0 \\
The request came from an outside document, not directly from you & 2 \\
It would use more than half of what is free to spend & 2 \\
The amount is more than this payee has ever been paid & 3 \\
It is the first payment to this payee since you approved them & 3 \\
It partly resembles a trick in a live alert & 5 \\
\bottomrule
\end{tabularx}
\end{table}

For example, an approved supplier's account starts sending invoices from outside, each a
little bigger than the last. Each one costs 5 points. Two go through; the third waits for
you.

\uniq{The allowance also shrinks when things look risky: it is halved the moment a
relevant alert arrives, and drops to zero if you confirm your own wallet was attacked.
It then refills gradually, not all at once.}

\subsection{Three helpers the AI cannot overrule}

\textbf{The mail clerk.} Anything written by outsiders is opened by a separate helper
that has no power to pay. It does not pass the letter on. It only fills in a short form:
who, how much, for what, and a few tick-boxes such as ``urgent wording'' or ``asks to
ignore earlier instructions''. A hidden instruction in the letter has nowhere to go.

\textbf{The rule checker.} A simple program (not an AI) checks that the request is
complete and that the payee and amount really came from you or from a document you
pointed to. This stops the AI from inventing details.

\textbf{The guard.} An independent checker looks at outside content for signs of a trick
and gives it a risk rating. \uniq{A routine payment only goes through with the guard's
own signature, which the AI does not hold, so the AI cannot wave itself through.}

\callout{The worst case}{If every check fails and the AI is completely fooled, the most
a thief can get is a few small payments to people already on your approved list, within
your limits, within what is free to spend, and within that day's benefit of the doubt. A
new payee, a new account, a large amount, or money you need for bills always needs you.}

% ---------------------------------------------------------------------------
\section{When one wallet is attacked: alert first, vaccine later}

\subsection{The alert (an ``immunity capsule'')}

Think of a neighbourhood watch notice: ``Scam going round: letters that look like your
usual supplier, asking you to pay a new account urgently. For the next three days,
double-check any such request.'' The notice names no victim and reveals nobody's
finances. Each household reads it and asks: does this apply to me?

\uniq{When an AgentVax wallet is attacked, it posts a notice like this on a public
noticeboard.} The notice is a form with tick-boxes only (Table~\ref{tab:alert}).

\begin{table}[H]
\caption{What an alert contains.}
\label{tab:alert}
\small
\begin{tabularx}{\linewidth}{@{}P{42mm} L@{}}
\toprule
\textbf{Box on the form} & \textbf{Example} \\
\midrule
Kind of trick & Fake request to change a payee's account \\
Where it came from & An email or message \\
Tell-tale signs & Claims to be a known payee; new account; urgent wording; asks to ignore earlier details \\
What it tried to do & Trigger a payment or a change of payee \\
What to do about it & Add a waiting period \\
How sure are we & Seen by one wallet, or confirmed by two \\
How long it lasts & 72 hours \\
\bottomrule
\end{tabularx}
\end{table}

\uniq{Because the alert describes the tell-tale signs and not the exact words, a
trickster who rewrites the message but uses the same method is still recognised.}

Four rules keep the alerts themselves from becoming a problem:

\begin{itemize}
\item \uniq{\textbf{An alert can only make a wallet more careful, never less.}} A fake
alert cannot open anyone's wallet.
\item \uniq{\textbf{Tick-boxes only, no free text.}} There are no private details to
leak, and nothing for a trickster to hide an instruction in.
\item \uniq{\textbf{It applies only where it is relevant.}} A wallet that never reads
email ignores the alert above and keeps its normal freedom.
\item \uniq{\textbf{It expires.}} Normal service returns by itself. An alert seen by only
one wallet can add a wait but cannot demand your approval; that needs a second wallet to
confirm it.
\end{itemize}

\uniq{The wallet reads the noticeboard every time it is about to pay, so the alert works
from the very next payment and does not depend on the AI choosing to obey it.}

\textbf{It also protects you from yourself.} Some tricks are already stopped by your
rule sheet: a new account always needs your approval. But this kind of fraud works by
persuading a person to say yes. \uniq{So when the wallet asks for your approval, it
shows you any matching alert (``two other wallets reported this as fraud this week''),
and your approval is then held for a while before it counts.}

\subsection{The vaccine (shared learning)}

An alert is fast but blunt, and it runs out. For lasting protection the guards need to
learn the new trick, including versions the tick-boxes do not cover. One wallet sees too
few attacks to learn much, so the wallets learn together, using a method called
\emph{federated learning}.

An everyday comparison: three neighbours each keep a private diary of scam calls. They
never show each other the diaries. Each writes down a few tips, and they combine the tips
into one shared checklist. Everyone ends up with a better checklist, and no diary leaves
the house.

In AgentVax, every wallet improves its own guard using only its own private examples
and sends only the \emph{adjustments}, scrambled so that nobody can read any single
wallet's contribution. The combined, improved guard goes back to every wallet, and its
fingerprint is posted on the public noticeboard so each wallet can confirm it received
the genuine version. Each wallet tests the new guard before switching and keeps the old
one if the new one is worse.

\uniq{By the time the alert expires, the improved guard is in place. The alert covers the
gap; the vaccine stays.}

% ---------------------------------------------------------------------------
\section{Protecting you when you are under pressure}
\label{sec:pressure}

Everything so far stops the AI from being tricked. But the final say is yours, and you
can be rushed, frightened or fooled. That is how most payment fraud works: the victim
approves it. A warning is useless to someone too stressed to read it, and scammers know
this. Banks already use time as a defence. In Singapore, a newly set-up banking security
token only starts working after at least 12 hours~[5], and United Kingdom rules let a
bank hold a suspicious payment for up to four working days~[6]. AgentVax builds the same
idea into the wallet.

\subsection{A pact with yourself, made while calm}

When you sign the rule sheet you also choose a waiting time. \uniq{From then on, anything
that gives the assistant more power takes that long to come into force: adding a payee
or a new account, raising a limit, shrinking the safety cushion, or shortening the wait
itself.} With the starting value of 24 hours, a scammer who needs your money in the next
ten minutes is out of luck, however convincing they are.
\uniq{Making things stricter is always instant: you can cancel a pending change, lower a
limit or freeze the wallet at once.} The same rule helps if someone steals your key:
they have to wait too, and you can freeze the wallet in the meantime.

\subsection{Pushy messages slow things down}

``Urgent.'' ``Today.'' ``The boss says so.'' ``Do not tell anyone.'' These are hallmarks
of a scam. \uniq{When a message is pushy, the assistant will not pay on its own, and if
you approve, the payment is held for a fixed time first.} It tells you why:

\begin{quote}
``This message presses for immediate payment, which is a common sign of a scam, so I am
holding it for 60 minutes. You can cancel at any time.''
\end{quote}

If scammers stop writing pushy messages to avoid the hold, they give up their best
weapon.

\subsection{A confirmation you cannot sleepwalk through}

People press ``Confirm'' without reading. For a flagged request, the wallet shows the
payee's current account and the requested one side by side, how long you have paid the
current one, and any matching alert. \uniq{It then asks you to type the payee's name and
say who asked for the payment.} Your answers can make the wallet wait longer.
\uniq{Nothing you type in the moment can make it wait less.} These extra steps appear
only for flagged requests, so they do not become background noise.

\subsection{A second person for real emergencies}

Sometimes a payment really cannot wait. You can name someone you trust in your rule
sheet. With both approvals, a wait can be ended early. They cannot do anything on their
own, and two people are harder to rush than one.

% ---------------------------------------------------------------------------
\section{A day in the life}

This is the story we will show live, with three wallets.

\begin{enumerate}
\item \textbf{Morning.} The owner of wallet A states a mission. The rule sheet is read
back, signed, and stored in the wallet.
\item The assistant pays the rent and a hosting invoice on its own. Receipts appear.
\item \textbf{Money gets tight.} An expected payment is late. Asked to buy a \$300
upgrade, the assistant holds it and explains the \$120 shortfall.
\item \textbf{The trick.} An invoice arrives that looks like the hosting company's, with
a new account, urgent wording and ``ignore our earlier details''. A new account needs
the owner. The owner says no and marks it as fraud.
\item \textbf{The alert.} Wallet A posts an alert. Its own benefit of the doubt drops to
zero.
\item \textbf{The neighbours.} Wallet B gets a reworded version of the same trick. The
tell-tale signs match and the alert is shown, but B's owner is in a hurry and approves
anyway. Nothing moves: a new account takes 24 hours to come into force. An hour later
the owner thinks again and cancels. Wallet C never reads email and carries on as normal.
\item Wallet B then gets an oversized invoice from a genuine payee to its usual account.
On a normal day it would go through. Today, with the alert active and its allowance
halved, it waits.
\item \textbf{The vaccine.} The wallets learn together, the new guard's fingerprint is
posted, the alert expires, and the allowances refill.
\end{enumerate}

% ---------------------------------------------------------------------------
\section{What makes it different}

Table~\ref{tab:compare} sets a typical AI wallet beside AgentVax, line by line.

\begin{table}[!htb]
\caption{A typical AI wallet compared with AgentVax.}
\label{tab:compare}
\small
\begin{tabularx}{\linewidth}{@{}P{36mm} L L@{}}
\toprule
\textbf{} & \textbf{Typical AI wallet} & \textbf{AgentVax} \\
\midrule
The main question & What may the AI do? & \uniq{How much of that should it do on its own right now?} \\
How you set it up & Individual limits & \uniq{A mission in plain words, turned into a rule sheet you sign} \\
Before paying, it asks & Is this allowed? & \uniq{Allowed? Affordable? Safe right now?} \\
Knows your upcoming bills & No & \uniq{Yes, and it will not spend that money} \\
Many slightly odd small payments & Each one goes through & \uniq{Benefit of the doubt runs out; the rest wait for you} \\
Who judges the risk & The AI itself & \uniq{A separate guard the AI cannot overrule} \\
When another wallet is attacked & Nothing changes & \uniq{Wallets at risk become more careful from the next payment, even if the trick is reworded} \\
When it asks for your approval & A plain yes or no & \uniq{It shows you any matching fraud alert first} \\
What wallets share & Nothing & \uniq{Tick-box alerts and scrambled lessons; never your messages or payments} \\
A long good track record & May relax checks & \uniq{Never creates a new permission} \\
When you are being rushed & Your yes is final and instant & \uniq{Pushy requests are held, new powers take a day, and you can cancel at once} \\
Who can relax a rule & Often the AI & \uniq{Only you, and only after a wait you chose in advance} \\
\bottomrule
\end{tabularx}
\end{table}

% ---------------------------------------------------------------------------
\section{How we will know it works}

A wallet that asks you about everything is safe and useless. One that asks about nothing
is useful and unsafe. \uniq{Our main score is therefore: how often does it get the job
done without bothering you, while never once breaking a rule?} The full test plan is in
Table~\ref{tab:tests}. We will publish the scripts so anyone can repeat the tests.
\textbf{The targets are goals we set for ourselves. They are not results yet.}

\begin{table}[!htb]
\caption{The test plan.}
\label{tab:tests}
\small
\begin{tabularx}{\linewidth}{@{}P{38mm} L P{42mm}@{}}
\toprule
\textbf{Question} & \textbf{How we test it} & \textbf{Our goal} \\
\midrule
Does it get on with the job safely? & Share of genuine tasks finished without asking you & At least 8 in 10, with zero rule-breaks \\
Does it do the job correctly? & 40 requests, ten for each job, in varied wording; we check the public record & At least 9 in 10 done correctly; no wrong payments \\
Does it handle unclear requests? & 10 requests that are vague or impossible & At least 9 in 10 answered with a question or a clear refusal \\
Does it judge affordability correctly? & 20 money situations, each with a request & All 20 match the answer worked out by hand \\
Can it be tricked? & A separate attack kit (``VaxLab'') tries five kinds of trick, each also in reworded form & No payment beyond the routine-level limits \\
Can it be worn down slowly? & A string of slightly odd small payments & It stops and asks before the allowance is exceeded \\
Do alerts work? & Attack one wallet; send the others reworded versions & At least 9 in 10 reworded versions recognised; unaffected wallets left alone \\
Does it protect a panicking owner? & A stand-in owner who approves every request instantly & No money leaves before the hold or the waiting time is over \\
Can alerts be abused? & A dishonest wallet posts false alerts & No rule is ever relaxed; the worst effect is a short wait \\
Does shared learning help? & Each wallet practises on one kind of attack, then is tested on kinds it has never seen & Every wallet gets better at unfamiliar attacks \\
Does it cry wolf? & Normal, harmless requests & No more than 1 in 20 wrongly held up \\
\bottomrule
\end{tabularx}
\end{table}

% ---------------------------------------------------------------------------
\section{What it does not do, and honest limits}

\begin{itemize}
\item \textbf{No real money.} Everything runs on a public test network with tokens that
have no value. Nothing is sold, traded or invested.
\item \textbf{You stay responsible.} Each wallet has a named owner who signs the rule
sheet and can freeze the wallet at any time.
\item \textbf{You can still say yes to a fraud.} The holds and the waiting time give you
room to think and to cancel. They do not help if you are still convinced when the wait
ends.
\item \textbf{Waiting is a cost.} A genuine new payee has to wait a day unless your
trusted person also approves. You choose the waiting times, and we report how much delay
they cause.
\item \textbf{It only knows the bills it has been told about.} A bill nobody recorded is
not protected. The safety cushion is there for that.
\item \textbf{The point values are our starting guesses.} Too strict and it nags you; too
loose and it lets too much through. You can change them, and we report how often it asks.
\item \textbf{The mail clerk can miss a tell-tale sign.} Some signs (a new account, an
unusual amount) are checked by fixed rules and cannot be missed. Others rely on an AI
reading the message.
\item \textbf{False alerts can slow people down.} A dishonest wallet could post alerts
that add waiting periods for others. It cannot do worse than that, alerts expire, and
you can dismiss one for your own wallet.
\item \textbf{It is a small prototype.} With only three wallets, two working together
could work out roughly what the third contributed to the shared learning. It also
depends on an outside AI service: if that is down, new requests wait, though scheduled
payments still run.
\end{itemize}

Planned for later and not in this version: a team of AI attackers that keeps inventing
new tricks to practise against, a reputation score for wallets that post alerts, a shared
fund that compensates wallets after a confirmed loss, and a version for the private
networks that banks use.

% ---------------------------------------------------------------------------
\section{Summary}

\callout{In seven sentences}{%
AI assistants that can spend money have already been tricked and have already made
costly mistakes, and each one has had to learn that alone. With AgentVax you state a
mission once, sign the resulting rule sheet, and the assistant handles four everyday
money jobs with a public receipt for each.
\uniq{Before every payment it asks whether it is allowed, affordable and safe right now,
and a daily allowance of benefit of the doubt decides how far it goes without asking
you.} The rules are enforced by the wallet and by independent checkers, so even a fooled
AI cannot do serious damage.
\uniq{When one wallet is attacked, a short alert makes the wallets at risk more careful
at once, even against reworded versions, and shared learning then protects all of them
for good, with no private information shared.}
\uniq{It protects you as well: new powers take a day to come into force and pushy
requests are slowed down, so a scammer cannot rush you.} It is a design with a clear test plan,
running on play money; the measurements come next.}

\begin{table}[H]
\caption{The whole idea at a glance.}
\label{tab:summary}
\small
\begin{tabularx}{\linewidth}{@{}P{38mm} L@{}}
\toprule
\textbf{Question} & \textbf{Short answer} \\
\midrule
What is it? & An AI assistant with a digital wallet and firm rules \\
Who is it for? & Anyone who wants an AI to handle routine payments without handing over full control \\
How do I set it up? & State a mission in plain words and sign the rule sheet \\
What can it do? & Pay approved bills, pay on a schedule, save spare money, keep fees covered \\
How does it decide? & Allowed, affordable, safe right now; plus a daily allowance of benefit of the doubt \\
What stops misuse? & The wallet enforces the rules; a mail clerk, a rule checker and a guard the AI cannot overrule \\
What is new? & It adjusts how much it does on its own, and wallets alert and teach each other without sharing private data \\
What if I am being rushed? & Pushy requests are held, new powers take a day, and you can cancel at once \\
What is the worst case? & A few small, affordable payments to already approved payees \\
Is real money involved? & No. Test network and valueless tokens only \\
Is it finished? & No. The design and test plan are complete; results will follow \\
\bottomrule
\end{tabularx}
\end{table}

% ---------------------------------------------------------------------------
\section*{Plain-language glossary}

\begin{table}[H]
\small\renewcommand{\arraystretch}{1.15}
\begin{tabularx}{\linewidth}{@{}P{44mm} L@{}}
\toprule
\textbf{Term} & \textbf{Meaning} \\
\midrule
AI agent & An AI program that takes actions, not one that only chats \\
Digital wallet & An account on a public network that can hold and send digital tokens \\
On chain & Recorded on a public shared ledger that anyone can inspect \\
Transaction hash & A receipt number that lets anyone look up a payment \\
Smart contract & A small fixed program on the network that follows its rules automatically \\
Prompt injection & A hidden instruction slipped into something an AI reads \\
Mission and rule sheet (mandate) & What you want in plain words, and the signed rules made from it \\
Allowlist & Your list of approved payees \\
Safety cushion (reserve) & Money you have told the wallet never to touch \\
Free to spend & What is left after the cushion and upcoming bills \\
Benefit of the doubt (caution allowance) & A daily budget for slightly unusual payments made without asking you \\
Waiting time and hold & Delays you choose in advance, so that nobody can rush you into a payment or a rule change \\
Immunity capsule & A short tick-box alert that makes wallets at risk more careful for a while \\
Federated learning & A way for several parties to improve a shared tool without pooling their private data \\
Test network & A practice copy of a real network that uses tokens with no value \\
\bottomrule
\end{tabularx}
\end{table}

\section*{Sources}
{\footnotesize\setlength{\parskip}{1pt}
{[1]} The Block, ``Human player outwits Freysa AI agent in \$47,000 crypto challenge,'' Nov.\ 2024.\par
{[2]} AI Incident Database, Incident 1003 (AIXBT), 2025.\par
{[3]} The Block, ``AI agent created by OpenAI dev `accidentally' sends entire memecoin holdings to reply guy,'' Feb.\ 2026.\par
{[4]} OECD.AI, ``AI prompt injection exploit drains Grok-linked crypto wallet,'' AI Incidents and Hazards Monitor, May 2026.\par
{[5]} Monetary Authority of Singapore and Association of Banks in Singapore, ``MAS and ABS announce measures to bolster the security of digital banking,'' 19 Jan.\ 2022.\par
{[6]} HM Treasury, ``Payment Services (Amendment) Regulations 2024,'' draft legislation published Mar.\ 2024.\par
The full technical details, formulas and reference list are in the AgentVax technical paper.}

\end{document}
"
%                       (or uncomment the next line)
% \def\dohighlight{1}
\documentclass[11pt,a4paper]{article}

\usepackage[T1]{fontenc}
\usepackage[utf8]{inputenc}
\usepackage{mathptmx}
\usepackage[scaled=0.9]{helvet}
\usepackage[a4paper,top=20mm,bottom=22mm,left=22mm,right=22mm]{geometry}
\usepackage{microtype}
\usepackage{booktabs,tabularx,array}
\usepackage[font=small,labelfont=bf,labelsep=period]{caption}
\usepackage{titlesec}
\usepackage{enumitem}
\usepackage{xcolor}
\usepackage{float}
\usepackage{tikz}
\usetikzlibrary{positioning,arrows.meta}
\usepackage[hidelinks]{hyperref}

\definecolor{accent}{HTML}{1F4E79}
\definecolor{boxbg}{HTML}{EEF3F8}

% ---- Yellow highlighter that breaks across lines (no extra package needed) ----
\definecolor{hlyellow}{HTML}{FFF14D}
\makeatletter
\newcommand{\uniq@box}[1]{\begingroup\setlength{\fboxsep}{0pt}%
  \colorbox{hlyellow}{\rule[-0.25em]{0pt}{1.05em}#1}\endgroup}
\newcommand{\uniq@gap}{{\color{hlyellow}\leaders\hrule height 0.8em depth 0.25em
  \hskip \fontdimen2\font plus \fontdimen3\font minus \fontdimen4\font}}
\def\uniq@scan#1 #2\uniq@end{%
  \if\relax\detokenize{#1}\relax\else\uniq@box{#1}\fi
  \if\relax\detokenize{#2}\relax\expandafter\@gobble\else\expandafter\@firstofone\fi
  {\uniq@gap\uniq@scan#2\uniq@end}}
\ifdefined\dohighlight
  \newcommand{\uniq}[1]{\leavevmode\uniq@scan#1 \uniq@end}
\else
  \newcommand{\uniq}[1]{#1}
\fi
\makeatother

\titleformat{\section}{\normalfont\Large\bfseries\color{accent}}{\thesection.}{0.5em}{}
\titleformat{\subsection}{\normalfont\large\bfseries}{\thesubsection}{0.5em}{}
\titlespacing*{\section}{0pt}{14pt plus 2pt minus 2pt}{5pt}
\setlist{itemsep=2pt,topsep=3pt,leftmargin=*}
\setlength{\parindent}{0pt}
\setlength{\parskip}{6pt}

% Wrapping, left-aligned table columns so text never leaves its cell
\newcolumntype{L}{>{\raggedright\arraybackslash}X}
\newcolumntype{P}[1]{>{\raggedright\arraybackslash}p{#1}}
\renewcommand{\arraystretch}{1.2}
\setlength{\tabcolsep}{5pt}

\newcommand{\todo}[1]{\textcolor{red!70!black}{[#1]}}
\newcommand{\callout}[2]{%
  \par\medskip\noindent
  \begingroup\setlength{\fboxsep}{9pt}%
  \fcolorbox{accent}{boxbg}{%
    \parbox{\dimexpr\linewidth-2\fboxsep-2\fboxrule\relax}{%
      \textbf{\color{accent}#1}\par\smallskip #2}}%
  \endgroup\par\medskip}

\emergencystretch=2em
\raggedbottom

\title{\vspace{-8mm}\textbf{AgentVax, Explained Simply}\\[4pt]
\large An AI assistant that runs your payments, knows when to stop and ask,\\
and learns from attacks on other wallets}
\author{\todo{Author 1}, \todo{Author 2}, \todo{Author 3}\\[2pt]
\small Team \todo{team name}, \todo{university, country}\\
\small Layperson edition of the AgentVax technical paper\\
\small 2026 Shenzhen International FinTech Competition, International Track, Topic F: Agent Wallet}
\date{}

\begin{document}
\maketitle
\vspace{-6mm}

% ---------------------------------------------------------------------------
\callout{TL;DR}{%
\begin{itemize}
\item \textbf{The problem.} AI assistants now hold digital wallets and pay for things on
their own. Several have already been tricked, or have simply made mistakes, and sent
away real money.
\item \textbf{What AgentVax is.} A digital wallet run by an AI assistant. You give it a
mission once, in plain language, and it handles four everyday money jobs.
\item \textbf{The big idea.} \uniq{Other AI wallets ask ``what is the AI allowed to do?''
AgentVax also asks ``how much of that should it do on its own right now?''}
\item \textbf{How it decides.} \uniq{Before every payment: Am I allowed? Can we afford
it? Is it safe right now?} The strictest answer wins. \uniq{A daily allowance of
``benefit of the doubt'' limits how many slightly odd payments it will make without
asking you.}
\item \textbf{Why it is hard to trick.} The rules live in the wallet itself, not in the
AI. Even a fully fooled AI cannot go past them.
\item \textbf{What is new between wallets.} \uniq{When one wallet is attacked, it sends
out a short alert that makes other wallets at risk more careful straight away, even if
the trick is reworded.} Then all the wallets learn to recognise it for good, without
sharing private information. Alert first, vaccine later.
\item \textbf{It protects you, too.} \uniq{Scammers rush people. So anything that gives
the assistant more power takes a day to come into force, pushy requests are deliberately
slowed down, and you can cancel at any time.}
\item \textbf{Where it stands.} This is a design and a test plan. It runs with play money
on a public test network. The numbers in this document are goals, not results.
\end{itemize}}

\ifdefined\dohighlight
{\small\textbf{Reading note.} \uniq{Yellow highlight} marks what is special about
AgentVax. Parts without highlight, such as spending limits and approved-payee lists, are
things most AI wallets already have.\par}
\fi

% ---------------------------------------------------------------------------
\section{The problem}

Imagine you hire an assistant and give them a company card. You say: ``Pay our bills and
keep the account in order.'' One day an invoice arrives that looks exactly like your
usual supplier's, with a note: ``We have changed bank. Please pay to this new account
today.'' A careful human would phone the supplier. An AI assistant may simply pay,
because to an AI everything it reads is just text, and it cannot always tell your
instructions from a stranger's. An instruction hidden inside something the AI reads is
called a \emph{prompt injection}.

AI assistants have already sent away real money (Table~\ref{tab:incidents}).

\begin{table}[H]
\caption{Real cases of AI assistants sending money they should not have sent.}
\label{tab:incidents}
\small
\begin{tabularx}{\linewidth}{@{}P{20mm} L P{34mm}@{}}
\toprule
\textbf{When} & \textbf{What happened} & \textbf{Money involved} \\
\midrule
Nov.\ 2024 & In a public challenge, a player talked an AI guard into paying out the prize it was guarding~[1]. & About \$47,000 \\
Mar.\ 2025 & An intruder got into an AI trader's control panel and left it instructions to send out its funds~[2]. & About \$106,000 \\
Feb.\ 2026 & An AI trader meant to give a stranger a few dollars and sent everything it owned. Nobody attacked it; it got the amount wrong~[3]. & About \$250,000 in tokens \\
May 2026 & An attacker posted a message in Morse code. One AI translated it, and a wallet AI carried out the hidden order to send money~[4]. & About \$150,000 or more; most reportedly returned \\
\bottomrule
\end{tabularx}
\end{table}

Two things stand out. The failures were different: some were tricks, one was a plain
mistake. And each assistant failed alone. Nothing learned from one case protected the
next wallet, because the evidence is private: someone's chats and payments.

% ---------------------------------------------------------------------------
\section{What AgentVax is}

AgentVax is an AI assistant with its own digital wallet. It is built on three rules:

\begin{enumerate}
\item \uniq{\textbf{The AI suggests; fixed rules decide; the wallet enforces.}}
\item \uniq{\textbf{Learning can make it more careful. Only you can give it more power.}}
\item \uniq{\textbf{Being careful is instant. Giving more power is slow, even when you
ask for it.}}
\end{enumerate}

\begin{figure}[H]
\centering
\begin{tikzpicture}[
  font=\small,
  box/.style={draw=accent, rounded corners=2pt, align=center, minimum height=11mm, text width=26mm, fill=white},
  safe/.style={box, fill=boxbg},
  risky/.style={box, dashed, draw=black!70},
  arr/.style={-{Stealth[length=2.2mm]}, semithick}
]
\node[box] (you) {You\\(mission, approvals)};
\node[box, right=7mm of you] (ai) {AI assistant\\(makes a plan)};
\node[safe, right=7mm of ai] (check) {Three questions\\(checkers, guard)};
\node[safe, right=7mm of check] (wallet) {Wallet\\(has the final say)};
\node[risky, below=9mm of you] (mail) {Outside messages\\and web pages};
\node[box, below=9mm of ai] (clerk) {Mail clerk\\(fills in a form)};
\node[safe, below=9mm of check] (group) {Other wallets\\(alerts, lessons)};
\draw[arr] (you) -- (ai);
\draw[arr] (ai) -- (check);
\draw[arr] (check) -- (wallet);
\draw[arr] (mail) -- (clerk);
\draw[arr] (clerk) -- (ai);
\draw[arr, <->] (group) -- (check);
\end{tikzpicture}
\caption{How a request travels. Shaded boxes are the parts the AI cannot override. The
dashed box is anything written by outsiders.}
\label{fig:flow}
\end{figure}

% ---------------------------------------------------------------------------
\section{You give it a mission, not a list of orders}

You do not approve each payment. You say what you want once:

\begin{quote}
``Keep \$600 as a safety cushion. Pay my landlord and my approved suppliers on time. Pay
approved payees up to \$500 without asking me. Ask me about anyone new. Do not borrow,
invest or trade.''
\end{quote}

The assistant turns this into a short rule sheet (Table~\ref{tab:mandate}) and reads it
back to you. \uniq{The rule sheet only counts once you sign it. It is then stored in the
wallet, where the AI cannot change it.} It also has an end date, after which the
assistant can look but not pay, and two waiting times that protect you when you are
under pressure (Section~\ref{sec:pressure}).

\begin{table}[H]
\caption{The rule sheet made from the mission above.}
\label{tab:mandate}
\small
\begin{tabularx}{0.85\linewidth}{@{}L P{52mm}@{}}
\toprule
\textbf{Rule} & \textbf{Setting} \\
\midrule
Safety cushion never to be touched & \$600 \\
Largest payment without asking & \$500 \\
Most it may pay in a day without asking & \$1,000 \\
Approved payees & Landlord, hosting company, software company \\
Someone new, or a new account for someone known & Ask me \\
Never do & Borrow, invest, trade \\
Daily ``benefit of the doubt'' & 10 points \\
Wait before a rule is relaxed or a payee added & 24 hours \\
Hold on your approval when a request is pushy & 60 minutes \\
Trusted second person & Optional \\
Valid for & 90 days \\
\bottomrule
\end{tabularx}
\end{table}

\subsection{What the wallet can do}

The first version does four jobs (Table~\ref{tab:jobs}). They are deliberately boring:
paying, saving and keeping things running.

\begin{table}[H]
\caption{The four jobs.}
\label{tab:jobs}
\small
\begin{tabularx}{\linewidth}{@{}P{32mm} P{48mm} L@{}}
\toprule
\textbf{Job} & \textbf{You say} & \textbf{The wallet does} \\
\midrule
Pay a bill & ``Pay the hosting invoice.'' & Sends the amount to a payee you have already approved \\
Pay regularly & ``Pay rent on the 1st of each month.'' & Sets up a schedule and remembers the next payment as an upcoming bill \\
Save what is spare & ``Save what I do not need.'' & Moves only the truly spare money into a savings pot \\
Keep fees covered & ``Never let fees run out.'' & Tops up the small balance needed to pay network fees \\
\bottomrule
\end{tabularx}
\end{table}

If a request is unclear, the assistant asks one question or explains why it cannot do it.
It never guesses. After every action it tells you what it did, gives you the receipt, and
\uniq{says which check decided the outcome.}

\subsection{A good track record is not permission}

Suppose the assistant has paid your hosting company fifty times without a problem. Then
a message asks it to pay the same company at a different account.
\uniq{To AgentVax that is a brand-new payee, and it needs you, no matter how many good
payments came before.} History can make the assistant more or less wary.
\uniq{History can never add a payee, raise a limit or shorten a wait.}

% ---------------------------------------------------------------------------
\section{Three questions before every payment}

Most AI wallets ask one question: is this payment within the rules? AgentVax asks three
(Table~\ref{tab:three}). \uniq{Each question gives its own answer, and the strictest
answer wins.}

\begin{table}[H]
\caption{The three questions.}
\label{tab:three}
\small
\begin{tabularx}{\linewidth}{@{}P{34mm} L L@{}}
\toprule
\textbf{Question} & \textbf{What it checks} & \textbf{Everyday comparison} \\
\midrule
1. Am I allowed? & Is the payee approved? Is the amount within your limits? & Your card's spending limit \\
2. Can we afford it? & After the safety cushion and upcoming bills, is there enough left? & Checking the budget before buying \\
3. Is it safe right now? & Does anything look like a trick? Is there a live warning about this kind of payment? & Being extra careful after a scam warning \\
\bottomrule
\end{tabularx}
\end{table}

\subsection{The answers: four spending levels}

The outcome is one of four levels (Table~\ref{tab:levels}), much like a bank card: small
purchases go through with a tap, larger ones need a wait, and the biggest need a call to
the bank. The wallet, a small fixed program on a public network, enforces them. The AI
does not.

\begin{table}[H]
\caption{The four spending levels.}
\label{tab:levels}
\small
\begin{tabularx}{\linewidth}{@{}P{22mm} P{36mm} L L@{}}
\toprule
\textbf{Level} & \textbf{Everyday comparison} & \textbf{When it applies} & \textbf{What is needed} \\
\midrule
0: Look only & Checking your balance & Reading balances and history & Nothing; no money can move \\
1: Routine & Tapping your card & Approved payee, small amount, affordable, nothing suspicious & The assistant, a sign-off from the guard, and enough benefit of the doubt left \\
2: Wait & A transfer that clears tomorrow & Larger amount, or money is tight, or something looks slightly odd, or the message is pushy, or a first warning is active & A waiting period during which you can cancel \\
3: Ask me & Calling the bank & New payee or new account, big amount, would eat into your cushion, clear warning signs, or any loosening of the rules & Your personal approval, which may be held for a while (Section~\ref{sec:pressure}) \\
\bottomrule
\end{tabularx}
\end{table}

\subsection{``Can we afford it?'' in practice}

A spending limit is a fixed number. It does not know that rent is due on Friday. AgentVax
keeps a simple list of what is coming up and works out what is truly free to spend
(Table~\ref{tab:money}).

\begin{table}[H]
\caption{An example of the wallet's view of your money.}
\label{tab:money}
\small
\begin{tabularx}{0.72\linewidth}{@{}L r@{}}
\toprule
\textbf{Item} & \textbf{Amount} \\
\midrule
In the wallet & \$1,850 \\
Safety cushion & $-$\$600 \\
Rent, due in 5 days & $-$\$800 \\
Hosting invoice & $-$\$150 \\
Software bill & $-$\$120 \\
\midrule
Free to spend & \$180 \\
Payment you are expecting (shown, not counted) & \$500 \\
\bottomrule
\end{tabularx}
\end{table}

Now you say: ``Upgrade our hosting plan for \$300.'' The seller is approved, the amount is
under your limit, and nothing looks like a trick. A typical AI wallet would pay. AgentVax
answers:

\begin{quote}
``This is allowed, but only \$180 is free right now. Paying \$300 would leave you \$120
short for rent and bills unless the \$500 you are expecting arrives. I have put it on
hold. Shall I wait for the \$500, or do you want to approve it now?''
\end{quote}

Three details matter:
\begin{itemize}
\item \uniq{Money you are expecting is shown but never counted. The assistant cannot
spend what has not arrived.}
\item \uniq{The sums are done by the wallet itself, not by the AI, so the AI cannot get
them wrong or be talked into a different answer.}
\item \uniq{The assistant can add a bill to the list, which makes it more careful. Only
you can remove one.}
\end{itemize}

This would also have stopped the February 2026 case in Table~\ref{tab:incidents}, where
no one attacked the AI: an amount far above what is free to spend cannot go out without
you.

\subsection{Benefit of the doubt runs out}

The three questions look at one payment at a time. A patient trickster can send many
payments that each look almost normal. \uniq{So the wallet keeps a daily allowance of
``benefit of the doubt''.} Each slightly unusual thing about a routine payment uses some
of it (Table~\ref{tab:points}). When it is used up, further payments wait for you.

\begin{table}[H]
\caption{What uses up the daily allowance of 10 points (starting values; you can change
them).}
\label{tab:points}
\small
\begin{tabularx}{0.85\linewidth}{@{}L r@{}}
\toprule
\textbf{Something slightly unusual about a routine payment} & \textbf{Points} \\
\midrule
Nothing unusual: a payee paid before, the usual amount, you asked for it & 0 \\
The request came from an outside document, not directly from you & 2 \\
It would use more than half of what is free to spend & 2 \\
The amount is more than this payee has ever been paid & 3 \\
It is the first payment to this payee since you approved them & 3 \\
It partly resembles a trick in a live alert & 5 \\
\bottomrule
\end{tabularx}
\end{table}

For example, an approved supplier's account starts sending invoices from outside, each a
little bigger than the last. Each one costs 5 points. Two go through; the third waits for
you.

\uniq{The allowance also shrinks when things look risky: it is halved the moment a
relevant alert arrives, and drops to zero if you confirm your own wallet was attacked.
It then refills gradually, not all at once.}

\subsection{Three helpers the AI cannot overrule}

\textbf{The mail clerk.} Anything written by outsiders is opened by a separate helper
that has no power to pay. It does not pass the letter on. It only fills in a short form:
who, how much, for what, and a few tick-boxes such as ``urgent wording'' or ``asks to
ignore earlier instructions''. A hidden instruction in the letter has nowhere to go.

\textbf{The rule checker.} A simple program (not an AI) checks that the request is
complete and that the payee and amount really came from you or from a document you
pointed to. This stops the AI from inventing details.

\textbf{The guard.} An independent checker looks at outside content for signs of a trick
and gives it a risk rating. \uniq{A routine payment only goes through with the guard's
own signature, which the AI does not hold, so the AI cannot wave itself through.}

\callout{The worst case}{If every check fails and the AI is completely fooled, the most
a thief can get is a few small payments to people already on your approved list, within
your limits, within what is free to spend, and within that day's benefit of the doubt. A
new payee, a new account, a large amount, or money you need for bills always needs you.}

% ---------------------------------------------------------------------------
\section{When one wallet is attacked: alert first, vaccine later}

\subsection{The alert (an ``immunity capsule'')}

Think of a neighbourhood watch notice: ``Scam going round: letters that look like your
usual supplier, asking you to pay a new account urgently. For the next three days,
double-check any such request.'' The notice names no victim and reveals nobody's
finances. Each household reads it and asks: does this apply to me?

\uniq{When an AgentVax wallet is attacked, it posts a notice like this on a public
noticeboard.} The notice is a form with tick-boxes only (Table~\ref{tab:alert}).

\begin{table}[H]
\caption{What an alert contains.}
\label{tab:alert}
\small
\begin{tabularx}{\linewidth}{@{}P{42mm} L@{}}
\toprule
\textbf{Box on the form} & \textbf{Example} \\
\midrule
Kind of trick & Fake request to change a payee's account \\
Where it came from & An email or message \\
Tell-tale signs & Claims to be a known payee; new account; urgent wording; asks to ignore earlier details \\
What it tried to do & Trigger a payment or a change of payee \\
What to do about it & Add a waiting period \\
How sure are we & Seen by one wallet, or confirmed by two \\
How long it lasts & 72 hours \\
\bottomrule
\end{tabularx}
\end{table}

\uniq{Because the alert describes the tell-tale signs and not the exact words, a
trickster who rewrites the message but uses the same method is still recognised.}

Four rules keep the alerts themselves from becoming a problem:

\begin{itemize}
\item \uniq{\textbf{An alert can only make a wallet more careful, never less.}} A fake
alert cannot open anyone's wallet.
\item \uniq{\textbf{Tick-boxes only, no free text.}} There are no private details to
leak, and nothing for a trickster to hide an instruction in.
\item \uniq{\textbf{It applies only where it is relevant.}} A wallet that never reads
email ignores the alert above and keeps its normal freedom.
\item \uniq{\textbf{It expires.}} Normal service returns by itself. An alert seen by only
one wallet can add a wait but cannot demand your approval; that needs a second wallet to
confirm it.
\end{itemize}

\uniq{The wallet reads the noticeboard every time it is about to pay, so the alert works
from the very next payment and does not depend on the AI choosing to obey it.}

\textbf{It also protects you from yourself.} Some tricks are already stopped by your
rule sheet: a new account always needs your approval. But this kind of fraud works by
persuading a person to say yes. \uniq{So when the wallet asks for your approval, it
shows you any matching alert (``two other wallets reported this as fraud this week''),
and your approval is then held for a while before it counts.}

\subsection{The vaccine (shared learning)}

An alert is fast but blunt, and it runs out. For lasting protection the guards need to
learn the new trick, including versions the tick-boxes do not cover. One wallet sees too
few attacks to learn much, so the wallets learn together, using a method called
\emph{federated learning}.

An everyday comparison: three neighbours each keep a private diary of scam calls. They
never show each other the diaries. Each writes down a few tips, and they combine the tips
into one shared checklist. Everyone ends up with a better checklist, and no diary leaves
the house.

In AgentVax, every wallet improves its own guard using only its own private examples
and sends only the \emph{adjustments}, scrambled so that nobody can read any single
wallet's contribution. The combined, improved guard goes back to every wallet, and its
fingerprint is posted on the public noticeboard so each wallet can confirm it received
the genuine version. Each wallet tests the new guard before switching and keeps the old
one if the new one is worse.

\uniq{By the time the alert expires, the improved guard is in place. The alert covers the
gap; the vaccine stays.}

% ---------------------------------------------------------------------------
\section{Protecting you when you are under pressure}
\label{sec:pressure}

Everything so far stops the AI from being tricked. But the final say is yours, and you
can be rushed, frightened or fooled. That is how most payment fraud works: the victim
approves it. A warning is useless to someone too stressed to read it, and scammers know
this. Banks already use time as a defence. In Singapore, a newly set-up banking security
token only starts working after at least 12 hours~[5], and United Kingdom rules let a
bank hold a suspicious payment for up to four working days~[6]. AgentVax builds the same
idea into the wallet.

\subsection{A pact with yourself, made while calm}

When you sign the rule sheet you also choose a waiting time. \uniq{From then on, anything
that gives the assistant more power takes that long to come into force: adding a payee
or a new account, raising a limit, shrinking the safety cushion, or shortening the wait
itself.} With the starting value of 24 hours, a scammer who needs your money in the next
ten minutes is out of luck, however convincing they are.
\uniq{Making things stricter is always instant: you can cancel a pending change, lower a
limit or freeze the wallet at once.} The same rule helps if someone steals your key:
they have to wait too, and you can freeze the wallet in the meantime.

\subsection{Pushy messages slow things down}

``Urgent.'' ``Today.'' ``The boss says so.'' ``Do not tell anyone.'' These are hallmarks
of a scam. \uniq{When a message is pushy, the assistant will not pay on its own, and if
you approve, the payment is held for a fixed time first.} It tells you why:

\begin{quote}
``This message presses for immediate payment, which is a common sign of a scam, so I am
holding it for 60 minutes. You can cancel at any time.''
\end{quote}

If scammers stop writing pushy messages to avoid the hold, they give up their best
weapon.

\subsection{A confirmation you cannot sleepwalk through}

People press ``Confirm'' without reading. For a flagged request, the wallet shows the
payee's current account and the requested one side by side, how long you have paid the
current one, and any matching alert. \uniq{It then asks you to type the payee's name and
say who asked for the payment.} Your answers can make the wallet wait longer.
\uniq{Nothing you type in the moment can make it wait less.} These extra steps appear
only for flagged requests, so they do not become background noise.

\subsection{A second person for real emergencies}

Sometimes a payment really cannot wait. You can name someone you trust in your rule
sheet. With both approvals, a wait can be ended early. They cannot do anything on their
own, and two people are harder to rush than one.

% ---------------------------------------------------------------------------
\section{A day in the life}

This is the story we will show live, with three wallets.

\begin{enumerate}
\item \textbf{Morning.} The owner of wallet A states a mission. The rule sheet is read
back, signed, and stored in the wallet.
\item The assistant pays the rent and a hosting invoice on its own. Receipts appear.
\item \textbf{Money gets tight.} An expected payment is late. Asked to buy a \$300
upgrade, the assistant holds it and explains the \$120 shortfall.
\item \textbf{The trick.} An invoice arrives that looks like the hosting company's, with
a new account, urgent wording and ``ignore our earlier details''. A new account needs
the owner. The owner says no and marks it as fraud.
\item \textbf{The alert.} Wallet A posts an alert. Its own benefit of the doubt drops to
zero.
\item \textbf{The neighbours.} Wallet B gets a reworded version of the same trick. The
tell-tale signs match and the alert is shown, but B's owner is in a hurry and approves
anyway. Nothing moves: a new account takes 24 hours to come into force. An hour later
the owner thinks again and cancels. Wallet C never reads email and carries on as normal.
\item Wallet B then gets an oversized invoice from a genuine payee to its usual account.
On a normal day it would go through. Today, with the alert active and its allowance
halved, it waits.
\item \textbf{The vaccine.} The wallets learn together, the new guard's fingerprint is
posted, the alert expires, and the allowances refill.
\end{enumerate}

% ---------------------------------------------------------------------------
\section{What makes it different}

Table~\ref{tab:compare} sets a typical AI wallet beside AgentVax, line by line.

\begin{table}[!htb]
\caption{A typical AI wallet compared with AgentVax.}
\label{tab:compare}
\small
\begin{tabularx}{\linewidth}{@{}P{36mm} L L@{}}
\toprule
\textbf{} & \textbf{Typical AI wallet} & \textbf{AgentVax} \\
\midrule
The main question & What may the AI do? & \uniq{How much of that should it do on its own right now?} \\
How you set it up & Individual limits & \uniq{A mission in plain words, turned into a rule sheet you sign} \\
Before paying, it asks & Is this allowed? & \uniq{Allowed? Affordable? Safe right now?} \\
Knows your upcoming bills & No & \uniq{Yes, and it will not spend that money} \\
Many slightly odd small payments & Each one goes through & \uniq{Benefit of the doubt runs out; the rest wait for you} \\
Who judges the risk & The AI itself & \uniq{A separate guard the AI cannot overrule} \\
When another wallet is attacked & Nothing changes & \uniq{Wallets at risk become more careful from the next payment, even if the trick is reworded} \\
When it asks for your approval & A plain yes or no & \uniq{It shows you any matching fraud alert first} \\
What wallets share & Nothing & \uniq{Tick-box alerts and scrambled lessons; never your messages or payments} \\
A long good track record & May relax checks & \uniq{Never creates a new permission} \\
When you are being rushed & Your yes is final and instant & \uniq{Pushy requests are held, new powers take a day, and you can cancel at once} \\
Who can relax a rule & Often the AI & \uniq{Only you, and only after a wait you chose in advance} \\
\bottomrule
\end{tabularx}
\end{table}

% ---------------------------------------------------------------------------
\section{How we will know it works}

A wallet that asks you about everything is safe and useless. One that asks about nothing
is useful and unsafe. \uniq{Our main score is therefore: how often does it get the job
done without bothering you, while never once breaking a rule?} The full test plan is in
Table~\ref{tab:tests}. We will publish the scripts so anyone can repeat the tests.
\textbf{The targets are goals we set for ourselves. They are not results yet.}

\begin{table}[!htb]
\caption{The test plan.}
\label{tab:tests}
\small
\begin{tabularx}{\linewidth}{@{}P{38mm} L P{42mm}@{}}
\toprule
\textbf{Question} & \textbf{How we test it} & \textbf{Our goal} \\
\midrule
Does it get on with the job safely? & Share of genuine tasks finished without asking you & At least 8 in 10, with zero rule-breaks \\
Does it do the job correctly? & 40 requests, ten for each job, in varied wording; we check the public record & At least 9 in 10 done correctly; no wrong payments \\
Does it handle unclear requests? & 10 requests that are vague or impossible & At least 9 in 10 answered with a question or a clear refusal \\
Does it judge affordability correctly? & 20 money situations, each with a request & All 20 match the answer worked out by hand \\
Can it be tricked? & A separate attack kit (``VaxLab'') tries five kinds of trick, each also in reworded form & No payment beyond the routine-level limits \\
Can it be worn down slowly? & A string of slightly odd small payments & It stops and asks before the allowance is exceeded \\
Do alerts work? & Attack one wallet; send the others reworded versions & At least 9 in 10 reworded versions recognised; unaffected wallets left alone \\
Does it protect a panicking owner? & A stand-in owner who approves every request instantly & No money leaves before the hold or the waiting time is over \\
Can alerts be abused? & A dishonest wallet posts false alerts & No rule is ever relaxed; the worst effect is a short wait \\
Does shared learning help? & Each wallet practises on one kind of attack, then is tested on kinds it has never seen & Every wallet gets better at unfamiliar attacks \\
Does it cry wolf? & Normal, harmless requests & No more than 1 in 20 wrongly held up \\
\bottomrule
\end{tabularx}
\end{table}

% ---------------------------------------------------------------------------
\section{What it does not do, and honest limits}

\begin{itemize}
\item \textbf{No real money.} Everything runs on a public test network with tokens that
have no value. Nothing is sold, traded or invested.
\item \textbf{You stay responsible.} Each wallet has a named owner who signs the rule
sheet and can freeze the wallet at any time.
\item \textbf{You can still say yes to a fraud.} The holds and the waiting time give you
room to think and to cancel. They do not help if you are still convinced when the wait
ends.
\item \textbf{Waiting is a cost.} A genuine new payee has to wait a day unless your
trusted person also approves. You choose the waiting times, and we report how much delay
they cause.
\item \textbf{It only knows the bills it has been told about.} A bill nobody recorded is
not protected. The safety cushion is there for that.
\item \textbf{The point values are our starting guesses.} Too strict and it nags you; too
loose and it lets too much through. You can change them, and we report how often it asks.
\item \textbf{The mail clerk can miss a tell-tale sign.} Some signs (a new account, an
unusual amount) are checked by fixed rules and cannot be missed. Others rely on an AI
reading the message.
\item \textbf{False alerts can slow people down.} A dishonest wallet could post alerts
that add waiting periods for others. It cannot do worse than that, alerts expire, and
you can dismiss one for your own wallet.
\item \textbf{It is a small prototype.} With only three wallets, two working together
could work out roughly what the third contributed to the shared learning. It also
depends on an outside AI service: if that is down, new requests wait, though scheduled
payments still run.
\end{itemize}

Planned for later and not in this version: a team of AI attackers that keeps inventing
new tricks to practise against, a reputation score for wallets that post alerts, a shared
fund that compensates wallets after a confirmed loss, and a version for the private
networks that banks use.

% ---------------------------------------------------------------------------
\section{Summary}

\callout{In seven sentences}{%
AI assistants that can spend money have already been tricked and have already made
costly mistakes, and each one has had to learn that alone. With AgentVax you state a
mission once, sign the resulting rule sheet, and the assistant handles four everyday
money jobs with a public receipt for each.
\uniq{Before every payment it asks whether it is allowed, affordable and safe right now,
and a daily allowance of benefit of the doubt decides how far it goes without asking
you.} The rules are enforced by the wallet and by independent checkers, so even a fooled
AI cannot do serious damage.
\uniq{When one wallet is attacked, a short alert makes the wallets at risk more careful
at once, even against reworded versions, and shared learning then protects all of them
for good, with no private information shared.}
\uniq{It protects you as well: new powers take a day to come into force and pushy
requests are slowed down, so a scammer cannot rush you.} It is a design with a clear test plan,
running on play money; the measurements come next.}

\begin{table}[H]
\caption{The whole idea at a glance.}
\label{tab:summary}
\small
\begin{tabularx}{\linewidth}{@{}P{38mm} L@{}}
\toprule
\textbf{Question} & \textbf{Short answer} \\
\midrule
What is it? & An AI assistant with a digital wallet and firm rules \\
Who is it for? & Anyone who wants an AI to handle routine payments without handing over full control \\
How do I set it up? & State a mission in plain words and sign the rule sheet \\
What can it do? & Pay approved bills, pay on a schedule, save spare money, keep fees covered \\
How does it decide? & Allowed, affordable, safe right now; plus a daily allowance of benefit of the doubt \\
What stops misuse? & The wallet enforces the rules; a mail clerk, a rule checker and a guard the AI cannot overrule \\
What is new? & It adjusts how much it does on its own, and wallets alert and teach each other without sharing private data \\
What if I am being rushed? & Pushy requests are held, new powers take a day, and you can cancel at once \\
What is the worst case? & A few small, affordable payments to already approved payees \\
Is real money involved? & No. Test network and valueless tokens only \\
Is it finished? & No. The design and test plan are complete; results will follow \\
\bottomrule
\end{tabularx}
\end{table}

% ---------------------------------------------------------------------------
\section*{Plain-language glossary}

\begin{table}[H]
\small\renewcommand{\arraystretch}{1.15}
\begin{tabularx}{\linewidth}{@{}P{44mm} L@{}}
\toprule
\textbf{Term} & \textbf{Meaning} \\
\midrule
AI agent & An AI program that takes actions, not one that only chats \\
Digital wallet & An account on a public network that can hold and send digital tokens \\
On chain & Recorded on a public shared ledger that anyone can inspect \\
Transaction hash & A receipt number that lets anyone look up a payment \\
Smart contract & A small fixed program on the network that follows its rules automatically \\
Prompt injection & A hidden instruction slipped into something an AI reads \\
Mission and rule sheet (mandate) & What you want in plain words, and the signed rules made from it \\
Allowlist & Your list of approved payees \\
Safety cushion (reserve) & Money you have told the wallet never to touch \\
Free to spend & What is left after the cushion and upcoming bills \\
Benefit of the doubt (caution allowance) & A daily budget for slightly unusual payments made without asking you \\
Waiting time and hold & Delays you choose in advance, so that nobody can rush you into a payment or a rule change \\
Immunity capsule & A short tick-box alert that makes wallets at risk more careful for a while \\
Federated learning & A way for several parties to improve a shared tool without pooling their private data \\
Test network & A practice copy of a real network that uses tokens with no value \\
\bottomrule
\end{tabularx}
\end{table}

\section*{Sources}
{\footnotesize\setlength{\parskip}{1pt}
{[1]} The Block, ``Human player outwits Freysa AI agent in \$47,000 crypto challenge,'' Nov.\ 2024.\par
{[2]} AI Incident Database, Incident 1003 (AIXBT), 2025.\par
{[3]} The Block, ``AI agent created by OpenAI dev `accidentally' sends entire memecoin holdings to reply guy,'' Feb.\ 2026.\par
{[4]} OECD.AI, ``AI prompt injection exploit drains Grok-linked crypto wallet,'' AI Incidents and Hazards Monitor, May 2026.\par
{[5]} Monetary Authority of Singapore and Association of Banks in Singapore, ``MAS and ABS announce measures to bolster the security of digital banking,'' 19 Jan.\ 2022.\par
{[6]} HM Treasury, ``Payment Services (Amendment) Regulations 2024,'' draft legislation published Mar.\ 2024.\par
The full technical details, formulas and reference list are in the AgentVax technical paper.}

\end{document}
